\documentclass[../main.tex]{subfiles}
\begin{document}

\chapter{スレッドと並行性}
\lessoninfo{
  video={P2L2},
  textbook={OSTEP \cite{ostep} ch.~26--32；Silberschatz ら \cite{silberschatz2018} ch.~4--5，8},
  keywords={スレッド，マルチスレッド，並行性，競合状態，ミューテックス，条件変数，リーダ・ライタ問題，スプリアスウェイクアップ，デッドロック，カーネルレベルスレッド，ユーザレベルスレッド，マルチスレッドモデル，ボス・ワーカーパターン，パイプラインパターン},
}

第2章で述べたプロセスは実行コンテキストを1つしか持たないので，一度に
1つの命令列しかたどれない．この章では，1つのアドレス空間を共有する複数の
実行コンテキストをプロセスに与えるスレッドと，そうしたスレッドが共有
データへのアクセスを協調させるための機構---ミューテックスと条件変数---を
導入する．説明は Birrell によるスレッドプログラミングのチュートリアル
\cite{birrell1989} に沿って進め，その後デッドロック，ユーザレベルスレッド
とカーネルレベルスレッドの関係，マルチスレッドアプリケーションを構成する
ための一般的なパターンを扱う．

\section{プロセスとスレッド}
\label{sec:l03-threads}

プロセスは，アドレス空間（ページテーブルに保持される仮想アドレスと物理
アドレスの対応付けを通じて到達するコード，データ，ヒープ），開いた
ファイルなどオペレーティングシステムがそのプロセスのために保持する資源，
そして実行コンテキスト，すなわちプログラムカウンタとスタックポインタを
含む CPU レジスタの値とスタックによって表される．\source{P2L2，プロセス
とスレッドの違い；Birrell \cite{birrell1989}；OSTEP ch.~26．}%
オペレーティングシステムはこれらすべてをプロセス制御ブロックに記録する．
このようなプロセスはシングルスレッドである．

\begin{definition}[スレッド]
  \term[すれっど]{スレッド}[thread]とは，プロセス内の1つの実行
  コンテキストである．各スレッドは独自のプログラムカウンタ，スタック
  ポインタ，レジスタ値，スタックを持つ．プロセスのすべてのスレッドは，
  そのプロセスのアドレス空間（コード，データ，ヒープ）と，開いたファイル
  などのその他の資源を共有する．
\end{definition}

1つのプロセスの中で複数のスレッドを実行することを
\term[まるちすれっど]{マルチスレッド}[multithreading]と呼び，そのような
プロセスを\emph{マルチスレッドプロセス}と呼ぶ．スレッドは仮想アドレスと
物理アドレスの対応付けを共有するので，すべてのスレッドが同じコードと
データにアクセスできる．しかし任意の時点で，各スレッドが実行している命令，
処理している入力の部分，進行中の手続き呼び出しの連鎖は，それぞれ異なりうる．
\cref{fig:l03-process-thread} のスレッドごとの状態は，まさにこれらの違い
を捉えている．プログラムカウンタとレジスタはスレッドがどこにいて何を計算
しているかを記録し，専用のスタックは実行中の呼び出しの局所変数と戻り
アドレスを保持する．

\begin{figure}[htbp]
  \centering
  % Single-threaded vs. multithreaded process: shared and per-thread state.
% Usage: % Single-threaded vs. multithreaded process: shared and per-thread state.
% Usage: % Single-threaded vs. multithreaded process: shared and per-thread state.
% Usage: \input{figures/l03-process-thread}   (width ~116 mm)
\begin{tikzpicture}[x=1mm, y=1mm, font=\footnotesize\sffamily,
  sh/.style={draw, fill=ln-box, minimum height=6mm, inner sep=0pt, anchor=south west},
  pt/.style={draw, fill=white, minimum height=6mm, minimum width=18mm,
             inner sep=0pt, anchor=south west},
  run/.style={-{Stealth[length=2mm]}, very thick, ln-accent},
  frame/.style={draw=ln-rule, rounded corners=2pt},
  ttl/.style={font=\footnotesize\sffamily\bfseries, anchor=south}]
  % ---- single-threaded process
  \draw[frame] (-2,-12) rectangle (44,22);
  \node[ttl] at (21,23) {\iflangja{シングルスレッドプロセス}{single-threaded process}};
  \node[sh, minimum width=14mm] at (0,14)  {\iflangja{コード}{code}};
  \node[sh, minimum width=14mm] at (14,14) {\iflangja{データ}{data}};
  \node[sh, minimum width=14mm] at (28,14) {\iflangja{ファイル}{files}};
  \node[pt] at (12,7) {\iflangja{レジスタ}{registers}};
  \node[pt] at (12,0) {\iflangja{スタック}{stack}};
  \draw[run] (21,-1.5) -- (21,-10);
  \node[anchor=west, text=ln-accent] at (22.5,-6) {\iflangja{スレッド}{thread}};
  % ---- multithreaded process
  \draw[frame] (52,-12) rectangle (116,22);
  \node[ttl] at (84,23) {\iflangja{マルチスレッドプロセス}{multithreaded process}};
  \node[sh, minimum width=20mm] at (54,14) {\iflangja{コード}{code}};
  \node[sh, minimum width=20mm] at (74,14) {\iflangja{データ}{data}};
  \node[sh, minimum width=20mm] at (94,14) {\iflangja{ファイル}{files}};
  \foreach \x in {54,75,96} {
    \node[pt] at (\x,7) {\iflangja{レジスタ}{registers}};
    \node[pt] at (\x,0) {\iflangja{スタック}{stack}};
    \draw[run] (\x+9,-1.5) -- (\x+9,-10);
  }
\end{tikzpicture}
   (width ~116 mm)
\begin{tikzpicture}[x=1mm, y=1mm, font=\footnotesize\sffamily,
  sh/.style={draw, fill=ln-box, minimum height=6mm, inner sep=0pt, anchor=south west},
  pt/.style={draw, fill=white, minimum height=6mm, minimum width=18mm,
             inner sep=0pt, anchor=south west},
  run/.style={-{Stealth[length=2mm]}, very thick, ln-accent},
  frame/.style={draw=ln-rule, rounded corners=2pt},
  ttl/.style={font=\footnotesize\sffamily\bfseries, anchor=south}]
  % ---- single-threaded process
  \draw[frame] (-2,-12) rectangle (44,22);
  \node[ttl] at (21,23) {\iflangja{シングルスレッドプロセス}{single-threaded process}};
  \node[sh, minimum width=14mm] at (0,14)  {\iflangja{コード}{code}};
  \node[sh, minimum width=14mm] at (14,14) {\iflangja{データ}{data}};
  \node[sh, minimum width=14mm] at (28,14) {\iflangja{ファイル}{files}};
  \node[pt] at (12,7) {\iflangja{レジスタ}{registers}};
  \node[pt] at (12,0) {\iflangja{スタック}{stack}};
  \draw[run] (21,-1.5) -- (21,-10);
  \node[anchor=west, text=ln-accent] at (22.5,-6) {\iflangja{スレッド}{thread}};
  % ---- multithreaded process
  \draw[frame] (52,-12) rectangle (116,22);
  \node[ttl] at (84,23) {\iflangja{マルチスレッドプロセス}{multithreaded process}};
  \node[sh, minimum width=20mm] at (54,14) {\iflangja{コード}{code}};
  \node[sh, minimum width=20mm] at (74,14) {\iflangja{データ}{data}};
  \node[sh, minimum width=20mm] at (94,14) {\iflangja{ファイル}{files}};
  \foreach \x in {54,75,96} {
    \node[pt] at (\x,7) {\iflangja{レジスタ}{registers}};
    \node[pt] at (\x,0) {\iflangja{スタック}{stack}};
    \draw[run] (\x+9,-1.5) -- (\x+9,-10);
  }
\end{tikzpicture}
   (width ~116 mm)
\begin{tikzpicture}[x=1mm, y=1mm, font=\footnotesize\sffamily,
  sh/.style={draw, fill=ln-box, minimum height=6mm, inner sep=0pt, anchor=south west},
  pt/.style={draw, fill=white, minimum height=6mm, minimum width=18mm,
             inner sep=0pt, anchor=south west},
  run/.style={-{Stealth[length=2mm]}, very thick, ln-accent},
  frame/.style={draw=ln-rule, rounded corners=2pt},
  ttl/.style={font=\footnotesize\sffamily\bfseries, anchor=south}]
  % ---- single-threaded process
  \draw[frame] (-2,-12) rectangle (44,22);
  \node[ttl] at (21,23) {\iflangja{シングルスレッドプロセス}{single-threaded process}};
  \node[sh, minimum width=14mm] at (0,14)  {\iflangja{コード}{code}};
  \node[sh, minimum width=14mm] at (14,14) {\iflangja{データ}{data}};
  \node[sh, minimum width=14mm] at (28,14) {\iflangja{ファイル}{files}};
  \node[pt] at (12,7) {\iflangja{レジスタ}{registers}};
  \node[pt] at (12,0) {\iflangja{スタック}{stack}};
  \draw[run] (21,-1.5) -- (21,-10);
  \node[anchor=west, text=ln-accent] at (22.5,-6) {\iflangja{スレッド}{thread}};
  % ---- multithreaded process
  \draw[frame] (52,-12) rectangle (116,22);
  \node[ttl] at (84,23) {\iflangja{マルチスレッドプロセス}{multithreaded process}};
  \node[sh, minimum width=20mm] at (54,14) {\iflangja{コード}{code}};
  \node[sh, minimum width=20mm] at (74,14) {\iflangja{データ}{data}};
  \node[sh, minimum width=20mm] at (94,14) {\iflangja{ファイル}{files}};
  \foreach \x in {54,75,96} {
    \node[pt] at (\x,7) {\iflangja{レジスタ}{registers}};
    \node[pt] at (\x,0) {\iflangja{スタック}{stack}};
    \draw[run] (\x+9,-1.5) -- (\x+9,-10);
  }
\end{tikzpicture}

  \caption{シングルスレッドプロセスとマルチスレッドプロセス．網掛けの状態
    （コード，データ，開いたファイル）はすべてのスレッドが共有する．各
    スレッドは PC と SP を含む独自のレジスタと，独自のスタックを持つ．}
  \label{fig:l03-process-thread}
\end{figure}

マルチスレッドプロセスに対するオペレーティングシステムの表現は，それに
応じてシングルスレッドプロセスの PCB より多くの情報を含む．すべてのスレッドが
共有する情報は1つだけ保持し，各スレッドについてはその実行コンテキストを
記録する別個のデータ構造を保持する．

\section{スレッドの有用性}
\label{sec:l03-benefits}

複数の CPU を持つマシンでは，スレッドは2通りの方法でアプリケーションを
高速化できる．\source{P2L2，スレッドはなぜ有用か，単一 CPU における
マルチスレッドの利点，アプリケーションと OS コードにとっての利点；OSTEP
ch.~26．}

\begin{description}
  \item[同じコード，異なるデータ] すべてのスレッドが同じコードを実行し，
    それぞれが入力の異なる部分を処理する．例えば行列積では，各スレッドが
    結果の行の異なる部分集合を計算でき，スレッドは異なる CPU 上で並列に
    実行される．1つのタスクの作業をこのように分割することを
    \term[へいれつか]{並列化}[parallelization]と呼ぶ．
  \item[異なるコード] スレッドごとにプログラムの異なる部分を実行する．
    あるスレッドが入力を処理し別のスレッドが画面を描画することもあれば，
    種類の異なる要求にそれぞれ専用のスレッドを割り当てることもある．
    スレッドを特定のタスクに専念させることを
    \term[とっか]{特化}[specialization]と呼ぶ．
\end{description}

特化にはさらに2つの利点がある．第一に，スレッドを異なる方法で管理できる．
例えば，より重要なタスクや顧客を扱うスレッドに高い優先度を与えられる．
第二に，性能はスレッドの状態のうちどれだけがプロセッサキャッシュにあるか
に依存するので，コードの小さな部分を繰り返し実行するスレッドは，すべてを
実行するスレッドよりもホットキャッシュ（第2章）で動作することが多い．

\subsection{スレッドとプロセスの比較}

同じ並列性は，シングルスレッドのプロセスを複数実行することでも得られる．
それでもマルチスレッドによる実装は，3つの点でより効率的である．
\begin{itemize}
  \item \textbf{メモリ}: スレッドは1つのアドレス空間を共有するが，
    プロセスはそれぞれ独自のアドレス空間と独自の状態のコピーを必要とする．
    マルチスレッドアプリケーションは必要なメモリが少なく，したがって
    スワッピングを必要とする可能性も低い．
  \item \textbf{コンテキストスイッチ}: 同じプロセスのスレッド間の切り替え
    ではアドレス空間を切り替えない．仮想アドレスと物理アドレスの対応付け
    はそのまま残るので，プロセスのコンテキストスイッチのうちコストの高い
    部分を避けられる．
  \item \textbf{通信}: スレッドは共有変数を通じてデータを交換する．
    プロセスにはプロセス間通信（IPC）が必要である．メッセージパッシング
    はやり取りのたびにカーネルに入り，共有メモリはまずオペレーティング
    システムによって設定されなければならないので，プロセス間の通信は
    よりコストが高い．
\end{itemize}

\begin{example}
  あるアプリケーションが，プロセスごとに別々に保持しなければならない
  \SI{400}{\mega\byte} のデータ（オペレーティングシステムがプロセス間で
  共有できるコードページはここでは無視する）を持ち，スタックその他の
  実行コンテキストごとの状態に約 \SI{1}{\mega\byte} を必要とするとする．
  それぞれがデータのコピーを持つ4つのプロセスとして実行すると約
  $4\times 401 = \SI{1604}{\mega\byte}$ を占め，1つのプロセスの4つの
  スレッドとして実行すると約 $400 + 4\times 1 = \SI{404}{\mega\byte}$
  を占める．
\end{example}

\subsection{単一 CPU 上のスレッド}

スレッドは，CPU よりスレッドの方が多い場合でも有用である．スレッドが
ディスク読み出しのようなブロックする要求を発行すると，要求が完了するまで
そのスレッドにはすることがなく，その間の時間 $t_{\mathrm{idle}}$ だけ CPU
はアイドルになってしまう．代わりに CPU は別のスレッドに切り替え，要求が
完了したら切り替えて戻ることができる（\cref{fig:l03-idle-time}）．この
とき別のスレッドは $t_{\mathrm{idle}} - t_{\mathrm{ctx}}$ の CPU 時間を
得る一方，ブロックしたスレッドは本来より $t_{\mathrm{ctx}}$ だけ遅れて
再開する．2回の切り替えのコストは合わせて $2\,t_{\mathrm{ctx}}$ である．
切り替えが有利になるのは，得られる時間が遅れを上回るとき，すなわち
\begin{equation}
  t_{\mathrm{idle}} \;>\; 2\, t_{\mathrm{ctx}}
  \label{eq:l03-idle}
\end{equation}
が成り立つときである．スレッド間のコンテキストスイッチはアドレス空間の
切り替えを避けるので，その $t_{\mathrm{ctx}}$ はプロセス間の切り替えより
小さく，\cref{eq:l03-idle} はより短いアイドル期間でも成り立つ．したがって
マルチスレッドにより，1つの CPU でもプログラムはアイドル時間を他の
スレッドの実行に活用できる．

\begin{figure}[htbp]
  \centering
  % Hiding I/O idle time on a single CPU by switching to another thread.
% Usage: % Hiding I/O idle time on a single CPU by switching to another thread.
% Usage: % Hiding I/O idle time on a single CPU by switching to another thread.
% Usage: \input{figures/l03-idle-time}   (width ~112 mm)
\begin{tikzpicture}[x=1mm, y=1mm, font=\footnotesize\sffamily,
  ev/.style={-{Stealth[length=1.6mm]}},
  dim/.style={{Bar[width=2mm]}-{Bar[width=2mm]}, thin},
  cs/.style={draw, fill=ln-accent!55}]
  \node[anchor=east] at (-1,3) {CPU};
  \draw[fill=ln-box] (0,0) rectangle (28,6);   \node at (14,3) {T1};
  \draw[cs]          (28,0) rectangle (33,6);
  \draw[fill=white]  (33,0) rectangle (83,6);  \node at (58,3) {T2};
  \draw[cs]          (83,0) rectangle (88,6);
  \draw[fill=ln-box] (88,0) rectangle (110,6); \node at (99,3) {T1};
  % events
  \draw[ev] (28,13) node[above] {\iflangja{T1 が I/O でブロック}{T1 blocks on I/O}} -- (28,7);
  \draw[ev] (83,13) node[above] {\iflangja{I/O 完了}{I/O completes}} -- (83,7);
  % dimensions
  \draw[dim] (28,-3) -- (33,-3);
  \draw[dim] (83,-3) -- (88,-3);
  \node[below] at (30.5,-3.5) {$t_{\mathrm{ctx}}$};
  \node[below] at (85.5,-3.5) {$t_{\mathrm{ctx}}$};
  \draw[dim] (28,-11) -- (83,-11);
  \node[below] at (55.5,-11.5) {$t_{\mathrm{idle}}$};
\end{tikzpicture}
   (width ~112 mm)
\begin{tikzpicture}[x=1mm, y=1mm, font=\footnotesize\sffamily,
  ev/.style={-{Stealth[length=1.6mm]}},
  dim/.style={{Bar[width=2mm]}-{Bar[width=2mm]}, thin},
  cs/.style={draw, fill=ln-accent!55}]
  \node[anchor=east] at (-1,3) {CPU};
  \draw[fill=ln-box] (0,0) rectangle (28,6);   \node at (14,3) {T1};
  \draw[cs]          (28,0) rectangle (33,6);
  \draw[fill=white]  (33,0) rectangle (83,6);  \node at (58,3) {T2};
  \draw[cs]          (83,0) rectangle (88,6);
  \draw[fill=ln-box] (88,0) rectangle (110,6); \node at (99,3) {T1};
  % events
  \draw[ev] (28,13) node[above] {\iflangja{T1 が I/O でブロック}{T1 blocks on I/O}} -- (28,7);
  \draw[ev] (83,13) node[above] {\iflangja{I/O 完了}{I/O completes}} -- (83,7);
  % dimensions
  \draw[dim] (28,-3) -- (33,-3);
  \draw[dim] (83,-3) -- (88,-3);
  \node[below] at (30.5,-3.5) {$t_{\mathrm{ctx}}$};
  \node[below] at (85.5,-3.5) {$t_{\mathrm{ctx}}$};
  \draw[dim] (28,-11) -- (83,-11);
  \node[below] at (55.5,-11.5) {$t_{\mathrm{idle}}$};
\end{tikzpicture}
   (width ~112 mm)
\begin{tikzpicture}[x=1mm, y=1mm, font=\footnotesize\sffamily,
  ev/.style={-{Stealth[length=1.6mm]}},
  dim/.style={{Bar[width=2mm]}-{Bar[width=2mm]}, thin},
  cs/.style={draw, fill=ln-accent!55}]
  \node[anchor=east] at (-1,3) {CPU};
  \draw[fill=ln-box] (0,0) rectangle (28,6);   \node at (14,3) {T1};
  \draw[cs]          (28,0) rectangle (33,6);
  \draw[fill=white]  (33,0) rectangle (83,6);  \node at (58,3) {T2};
  \draw[cs]          (83,0) rectangle (88,6);
  \draw[fill=ln-box] (88,0) rectangle (110,6); \node at (99,3) {T1};
  % events
  \draw[ev] (28,13) node[above] {\iflangja{T1 が I/O でブロック}{T1 blocks on I/O}} -- (28,7);
  \draw[ev] (83,13) node[above] {\iflangja{I/O 完了}{I/O completes}} -- (83,7);
  % dimensions
  \draw[dim] (28,-3) -- (33,-3);
  \draw[dim] (83,-3) -- (88,-3);
  \node[below] at (30.5,-3.5) {$t_{\mathrm{ctx}}$};
  \node[below] at (85.5,-3.5) {$t_{\mathrm{ctx}}$};
  \draw[dim] (28,-11) -- (83,-11);
  \node[below] at (55.5,-11.5) {$t_{\mathrm{idle}}$};
\end{tikzpicture}

  \caption{T1 は I/O でブロックする．CPU はアイドルになる代わりに T2 に
    切り替え，I/O の完了を契機に切り替えて戻る．切り替えのコストがアイドル
    時間より小さければ，切り替えは有利になる（\cref{eq:l03-idle}）．}
  \label{fig:l03-idle-time}
\end{figure}

\begin{example}
  スレッドのコンテキストスイッチに \SI{0.05}{\milli\second}，プロセスの
  コンテキストスイッチに \SI{0.5}{\milli\second} かかるとする．スレッド
  なら \SI{0.1}{\milli\second} より長いブロッキング待ちはすべて別の
  スレッドが利用できるが，プロセスでは \SI{1}{\milli\second} より長い
  待ちに限られる．\SI{4}{\milli\second} の待ちは，別のスレッドに
  $4 - 0.05 = \SI{3.95}{\milli\second}$ の CPU 時間を与え，ブロックした
  スレッドを \SI{0.05}{\milli\second} 遅らせる．別のプロセスの場合，
  これらの値は \SI{3.5}{\milli\second} と \SI{0.5}{\milli\second} である．
\end{example}

\subsection{マルチスレッドのアプリケーションとカーネル}

アプリケーションもオペレーティングシステムも恩恵を受ける．マルチスレッド
アプリケーションは，複数の CPU では並列化による高速化を達成し，1つの CPU
ではアイドル時間を有効に活用する．マルチスレッドのオペレーティングシステムカーネルは，
一部のスレッドをアプリケーションのために実行し，他のスレッドをデーモンや
デバイスドライバなどのオペレーティングシステムサービスを提供するために
実行する．マルチプロセッサでは，これらのカーネルスレッドは異なる CPU 上で
並行に実行される．

\begin{keypoint}
  スレッドはプロセスのアドレス空間を共有するので，メモリ，コンテキスト
  スイッチ時間，通信の点でプロセスより安価である．同じ共有のために
  スレッドは互いに干渉しうるので，協調のための明示的な機構が必要となる．
\end{keypoint}

\section{スレッドの基本機構}
\label{sec:l03-mechanisms}

スレッドを実現するには3つのものが必要である．すなわち，スレッドを識別し
その資源使用を追跡できるようにするための，スレッドを表すデータ構造，
スレッドを生成・管理する機構，そして同じアドレス空間で同時に実行される
スレッドが安全に協調するための機構である．\source{P2L2，スレッドの基本
機構，スレッドの生成；Birrell \cite{birrell1989}；OSTEP ch.~26--27．}

2つのスレッドの実行が時間的に重なるとき，それらは並行に実行されるという．
1つの CPU 上で交互に実行される場合も，異なる CPU 上で同時に実行される場合
も同様である．こうして生じる\term[へいこうせい]{並行性}[concurrency]は，
プロセスにはない問題を生む．オペレーティングシステムは異なるプロセスに異なる対応付け
を与えるので，あるプロセスが別のプロセスのメモリに触れることはできない．
スレッドは対応付けを共有するので，どのスレッドもプロセスのどのデータでも
正当に読み書きできる．2つのスレッドが同じデータを同時に更新したり，一方
が読んでいる間に他方が変更したりすると，データは不整合な状態になりうる．

\begin{definition}[競合状態]
  \term[きょうごうじょうたい]{競合状態}[race condition]とは，計算結果が，
  共有データにアクセスするスレッド（少なくとも1つはそれを変更する）の
  相対的なタイミングに依存する状況である．同じデータに対する，
  同期されていない2つの並行アクセスで，少なくとも一方が書き込みである
  ものを\emph{データ競合}と呼ぶ．データ競合は競合状態のよくある原因である．
\end{definition}

スレッドが共有データへのアクセスと互いの進行を協調させるために用いる機構
を\term[どうき]{同期}[synchronization]機構と呼ぶ．この章では，どちらも
Birrell が述べている次の2つの基本機構を全体を通じて用いる．
\begin{itemize}
  \item \textbf{相互排除}: 一度に1つのスレッドだけがある操作を実行できる．
    これは\emph{ミューテックス}が提供する（\cref{sec:l03-mutex}）．
  \item \textbf{待機}: あるスレッドが，別のスレッドが特定の条件を成立
    させるまで待ち，そのスレッドによって起こされる．これは
    \emph{条件変数}が提供する（\cref{sec:l03-condvar}）．
\end{itemize}

\subsection{スレッドの生成}
\label{sec:l03-creation}

スレッドは，その識別子，プログラムカウンタ，スタックポインタその他の
レジスタ値，スタック，そしてスレッドライブラリがスレッドを管理するために
使う属性を含むスレッドデータ構造で表される．Birrell の記法では，
\term[すれっどせいせい]{スレッド生成}[thread creation]は呼び出し
\texttt{t1 = Fork(proc, args)} で行う．この呼び出しは新しいスレッドを
生成し，そのデータ構造を \texttt{t1} に格納し，新しいスレッドに
\texttt{proc(args)} の実行を開始させる．生成したスレッドは \texttt{Fork}
の次の文から実行を続けるので，この時点から両方のスレッドが並行に実行
される．名前に反して \texttt{Fork} は同じプロセス内にスレッドを生成する
ものであり，新しいプロセスを生成する第2章の Unix の \texttt{fork} とは
無関係である．

子スレッドが終了したら，その結果を親に届けなければならない．子が共有
アドレス空間内の決められた位置に結果を格納するか，親が
\term[じょいん]{ジョイン}[join]操作で結果を得る．\texttt{Join(t1)} は，
スレッド \texttt{t1} が完了するまで呼び出し元のスレッドをブロックし，
その結果を返す．この時点で子に割り当てられていた資源は解放される．これを
除けば親と子は対等である．どちらもプロセスのすべての資源にアクセスでき，
異なるのは実行コンテキストだけである．

\needspace{9\baselineskip}
\begin{example}
  親スレッドと子スレッドがどちらも共有リストに挿入する．
  \begin{lstlisting}[language=C, numbers=none]
Thread t1;
Shared_list list;

t1 = Fork(safe_insert, 7);  // 子が 7 を挿入
safe_insert(9);             // 親が 9 を挿入
Join(t1);                   // ここでは省略可
\end{lstlisting}
  どちらの挿入が先に起こるかはスレッドのスケジューリングに依存するので，
  先頭に挿入する場合，リストは最終的に $7\to 9$ または $9\to 7$ の
  どちらかになる．親は子の結果を使わないので，\texttt{Join} は正しさの
  ためには必要ない．
\end{example}

この例には問題が隠れている．リストの先頭に要素 $e$ を挿入するには，$e$ を
生成する，先頭ポインタを読んで $e$\texttt{.next} に格納する，$e$ の
アドレスを先頭ポインタに書き込む，という3つのステップが必要である．2つの
スレッドがこれらのステップを交互に実行すると，\cref{tab:l03-race} の
トレースが示すように挿入の1つが失われうる．

\begin{table}[htbp]
  \centering
  \caption{初めに要素 3 を1つだけ含むリストへの，同期されていない2つの
    挿入．どちらのスレッドも，いずれかが先頭を書き込む前に古い先頭を読む
    ので，ステップ~6 の書き込みがステップ~5 の書き込みを上書きし，
    要素~7 が失われる．}
  \label{tab:l03-race}
  \small
  \begin{tabular}{@{}clll@{}}
    \toprule
    ステップ & スレッド T1（7 を挿入） & スレッド T2（9 を挿入） & リスト \\
    \midrule
    1 & \texttt{e7} を生成            &                           & $3$ \\
    2 & \texttt{e7.next = head} ($\to 3$) &                      & $3$ \\
    3 &                              & \texttt{e9} を生成        & $3$ \\
    4 &                              & \texttt{e9.next = head} ($\to 3$) & $3$ \\
    5 & \texttt{head = e7}           &                           & $7\to 3$ \\
    6 &                              & \texttt{head = e9}        & $9\to 3$ \\
    \bottomrule
  \end{tabular}
\end{table}

\section{相互排除}
\label{sec:l03-mutex}

\cref{tab:l03-race} の競合は，挿入の3つのステップを一度に1つのスレッド
しか実行できなければ消える．\source{P2L2，ミューテックス；Birrell
\cite{birrell1989}；OSTEP ch.~28．}このとき3つのステップは，他のスレッド
から見て1つの分割不可能な操作のように振る舞わなければならない．すなわち，
ステップの間に他のスレッドがリストを観測したり変更したりすることはできない．
この要件を\term[そうごはいじょ]{相互排除}[mutual exclusion]と呼び，それを
強制する基本機構がミューテックスである．

\begin{definition}[ミューテックス]
  \term[みゅーてっくす]{ミューテックス}[mutex]（相互排除ロック）は，
  ロックとアンロックの2つの操作を持つ同期プリミティブである．空いている
  ミューテックスをロックしたスレッドはそれを獲得し，その所有者となる．
  別のスレッドが所有するミューテックスをロックしようとしたスレッドは，
  ミューテックスが解放されるまでブロックする．
\end{definition}

ミューテックスは，少なくともロックされているかどうか，どのスレッドが所有
しているか，そのミューテックスでブロックしているスレッドのリストを記録する
データ構造で表される．ミューテックスを保持している間に実行されるコードが
\term[くりてぃかるせくしょん]{クリティカルセクション}[critical section]%
である．これは，一度に1つのスレッドしか実行してはならない操作，典型的には
共有状態の更新を行うプログラムの部分である．

\needspace{6\baselineskip}
Birrell の記法では，クリティカルセクションをブロックとして書く．
\begin{lstlisting}[language=C, numbers=none]
Lock(m) {
  // クリティカルセクション
} // unlock
\end{lstlisting}
ミューテックスはブロックの終わりで暗黙に解放される．第4章の POSIX
スレッドを含む多くのスレッドインタフェースでは，代わりにクリティカル
セクションの前後で別々の呼び出し \texttt{lock(m)} と \texttt{unlock(m)}
を行い，クリティカルセクションを通るすべての経路がアンロックで終わるように
することはプログラマの責任となる．

\needspace{12\baselineskip}
\begin{example}
  リストを保護するミューテックスを用いると，\cref{sec:l03-creation}の
  挿入は安全になる．
  \begin{lstlisting}[language=C, numbers=none]
list<int> my_list;
Mutex m;

void safe_insert(int i) {
  Lock(m) {
    my_list.insert(i);
  } // unlock
}
\end{lstlisting}
  親が先に \texttt{m} をロックすると，子は親の挿入が完了するまで
  \texttt{Lock(m)} でブロックし，その後自分の挿入を行う．
  \cref{tab:l03-race} の状況では，リストは最終的に $7\to 9\to 3$ または
  $9\to 7\to 3$ となり，要素が失われることはない．
\end{example}

複数のスレッドがブロックしているミューテックスを所有者が解放すると，その
うち1つだけがミューテックスを獲得し，残りはブロックしたままである．
ミューテックスの仕様は，どれが獲得するかを定めていない．スレッドが要求した
順にミューテックスを獲得できる保証はなく，ミューテックスが解放された瞬間に
\texttt{Lock} を呼んだスレッドが，しばらく待っていたスレッドより先に獲得
することもある．したがって，プログラムの正しさは，待っているスレッドが
ミューテックスを獲得する順序に依存してはならない．

\section{条件変数}
\label{sec:l03-condvar}

他のスレッドが成立させる条件をスレッドが待たなければならない場合，相互排除
だけでは不十分である．\source{P2L2，生産者/消費者の例，条件変数，条件変数
の API；Birrell \cite{birrell1989}；OSTEP ch.~30．}ミューテックスを保持
しているスレッドは，単に停止して条件を待つことはできない．ミューテックス
を保持している間は，他のスレッドが共有状態を変更できないからである．
標準的な例は
\term[せいさんしゃしょうひしゃもんだい]{生産者・消費者問題}[producer-consumer problem]%
である．複数の生産者スレッドが共有リストに項目を挿入し，消費者スレッドは
リストが満杯になったらすべての項目を表示して取り除かなければならない．

ミューテックスだけを使う場合，消費者はミューテックスをロックし，リストが
満杯かどうかを調べ，満杯でなければミューテックスを解放して後で再試行
しなければならない．ほとんどの検査ではリストはまだ満杯でないので，この
繰り返しの検査は CPU 時間を浪費する．消費者は眠って待ち，リストを満杯に
する挿入を行った生産者によって起こされる方がよい．

\begin{definition}[条件変数]
  \term[じょうけんへんすう]{条件変数}[condition variable]は，ミューテックス
  と組み合わせて用いる同期プリミティブである．スレッドは共有状態に関する
  ある条件が成り立つまで条件変数上で待ち，他のスレッドは条件が成り立った
  可能性があるときに条件変数を通じて待っているスレッドに通知する．
\end{definition}

\needspace{8\baselineskip}
条件変数 \texttt{list\_full} を用いると，両者は次のようになる．
\begin{lstlisting}[language=C, numbers=none]
// 消費者: print_and_clear
Lock(m) {
  while (my_list.not_full())
    Wait(m, list_full);
  my_list.print_and_remove_all();
} // unlock
\end{lstlisting}
\needspace{7\baselineskip}
\begin{lstlisting}[language=C, numbers=none]
// 生産者: safe_insert
Lock(m) {
  my_list.insert(my_thread_id);
  if (my_list.full())
    Signal(list_full);
} // unlock
\end{lstlisting}

\subsection{条件変数のインタフェース}

Birrell のインタフェースには，条件変数 \texttt{cond} に対する3つの操作が
ある．
\begin{description}
  \item[\texttt{Wait(mutex, cond)}] \term[waitそうさ]{wait操作}[wait operation]%
    は \texttt{mutex} を保持した状態で呼び出す．この操作はミューテックスを
    解放し，スレッドを \texttt{cond} の待ちキューで停止させる．スレッドが
    起こされると，操作は戻る前にミューテックスを再獲得する．
  \item[\texttt{Signal(cond)}] \term[signalそうさ]{signal操作}[signal operation]%
    は，\texttt{cond} で待っているスレッドがあれば，そのうち1つを起こす．
  \item[\texttt{Broadcast(cond)}]
    \term[broadcastそうさ]{broadcast操作}[broadcast operation]は，
    \texttt{cond} で待っているすべてのスレッドを起こす．
\end{description}

\needspace{11\baselineskip}
条件変数は，そのミューテックスへの参照と待っているスレッドのリストを保持
するデータ構造で表される．wait 操作は次のように進む．
\begin{lstlisting}[language=C, numbers=none]
Wait(mutex, cond) {
  // アトミックに mutex を解放し，
  //   スレッドを cond の待ちキューに入れる
  // ... 待つ ...
  // スレッドを待ちキューから取り除く
  // mutex を再獲得する
  // 戻る
}
\end{lstlisting}
待っている間にミューテックスを解放することは不可欠である．さもなければ，
他のどのスレッドもクリティカルセクションに入って状態を変更し，signal する
ことができない．戻る前にミューテックスを再獲得することで，待っていた
スレッドが共有状態を再検査するときに再びミューテックスを保持していること
が保証される．その結果，起こされたスレッドはミューテックスを再び保持する
まで先に進めない．ミューテックスがまだ保持されている場合，例えば signal
したスレッドが保持している場合には，起こされたスレッドは条件変数の待ち
キューからミューテックスの待ちキューへ移る（\cref{fig:l03-cv-queues}）．
broadcast の後は起こされたすべてのスレッドがミューテックスを奪い合うので，
それらは1つずつ wait 操作から戻る．\margin{条件変数の signal 操作は，
第1章の Unix シグナルとは無関係である．}

\begin{figure}[htbp]
  \centering
  % Life of a thread around a mutex m and a condition variable c.
% Usage: % Life of a thread around a mutex m and a condition variable c.
% Usage: % Life of a thread around a mutex m and a condition variable c.
% Usage: \input{figures/l03-cv-queues}   (width ~116 mm)
\begin{tikzpicture}[x=1mm, y=1mm, font=\footnotesize\sffamily,
  st/.style={draw, rounded corners=2pt, fill=ln-box, align=center,
             minimum height=11mm, inner sep=2pt},
  >={Stealth[length=1.8mm]}, lab/.style={font=\scriptsize\sffamily, align=center}]
  \node[st, minimum width=26mm, text width=24mm] (mq) at (13,0)
    {\iflangja{m の待ちキュー}{wait queue of mutex m}};
  \node[st, minimum width=28mm, text width=26mm, fill=white, thick] (cs) at (58,0)
    {\iflangja{クリティカルセクション\\（m を保持）}{critical section\\(owns m)}};
  \node[st, minimum width=26mm, text width=24mm] (cq) at (103,0)
    {\iflangja{c の待ちキュー}{wait queue of condition c}};
  \node[font=\footnotesize\ttfamily, anchor=south] (lk) at (35,21) {Lock(m)};
  \draw (lk.south) -- (35,16);
  \draw[->] (35,16) -| node[lab, pos=0.4, above] {\iflangja{m は使用中}{m held}} (mq.north);
  \draw[->] (35,16) -| node[lab, pos=0.5, above] {\iflangja{m は空き}{m free}} (cs.north);
  \draw[->] (mq) -- node[lab, above] {\iflangja{m を獲得}{acquire m}} (cs);
  \draw[->] (cs) -- node[lab, above] {\texttt{Wait(m,c)}}
                    node[lab, below] {\iflangja{m を解放}{releases m}} (cq);
  \draw[->] ([xshift=8mm]cs.north) |- node[lab, pos=0.75, above] {\texttt{Unlock(m)}} (94,16);
  % signalled: re-acquire m directly if it is free, otherwise queue on m
  \draw (cq.south) -- node[lab, left, pos=0.55] {\texttt{Signal(c)} / \texttt{Broadcast(c)}} (103,-15);
  \draw[->] (103,-15) -| node[lab, pos=0.75, left] {\iflangja{m は空き}{m free}} ([xshift=8mm]cs.south);
  \draw[->] (103,-15) -- (103,-21) -| node[lab, pos=0.25, below] {\iflangja{m は使用中}{m held}} (mq.south);
\end{tikzpicture}
   (width ~116 mm)
\begin{tikzpicture}[x=1mm, y=1mm, font=\footnotesize\sffamily,
  st/.style={draw, rounded corners=2pt, fill=ln-box, align=center,
             minimum height=11mm, inner sep=2pt},
  >={Stealth[length=1.8mm]}, lab/.style={font=\scriptsize\sffamily, align=center}]
  \node[st, minimum width=26mm, text width=24mm] (mq) at (13,0)
    {\iflangja{m の待ちキュー}{wait queue of mutex m}};
  \node[st, minimum width=28mm, text width=26mm, fill=white, thick] (cs) at (58,0)
    {\iflangja{クリティカルセクション\\（m を保持）}{critical section\\(owns m)}};
  \node[st, minimum width=26mm, text width=24mm] (cq) at (103,0)
    {\iflangja{c の待ちキュー}{wait queue of condition c}};
  \node[font=\footnotesize\ttfamily, anchor=south] (lk) at (35,21) {Lock(m)};
  \draw (lk.south) -- (35,16);
  \draw[->] (35,16) -| node[lab, pos=0.4, above] {\iflangja{m は使用中}{m held}} (mq.north);
  \draw[->] (35,16) -| node[lab, pos=0.5, above] {\iflangja{m は空き}{m free}} (cs.north);
  \draw[->] (mq) -- node[lab, above] {\iflangja{m を獲得}{acquire m}} (cs);
  \draw[->] (cs) -- node[lab, above] {\texttt{Wait(m,c)}}
                    node[lab, below] {\iflangja{m を解放}{releases m}} (cq);
  \draw[->] ([xshift=8mm]cs.north) |- node[lab, pos=0.75, above] {\texttt{Unlock(m)}} (94,16);
  % signalled: re-acquire m directly if it is free, otherwise queue on m
  \draw (cq.south) -- node[lab, left, pos=0.55] {\texttt{Signal(c)} / \texttt{Broadcast(c)}} (103,-15);
  \draw[->] (103,-15) -| node[lab, pos=0.75, left] {\iflangja{m は空き}{m free}} ([xshift=8mm]cs.south);
  \draw[->] (103,-15) -- (103,-21) -| node[lab, pos=0.25, below] {\iflangja{m は使用中}{m held}} (mq.south);
\end{tikzpicture}
   (width ~116 mm)
\begin{tikzpicture}[x=1mm, y=1mm, font=\footnotesize\sffamily,
  st/.style={draw, rounded corners=2pt, fill=ln-box, align=center,
             minimum height=11mm, inner sep=2pt},
  >={Stealth[length=1.8mm]}, lab/.style={font=\scriptsize\sffamily, align=center}]
  \node[st, minimum width=26mm, text width=24mm] (mq) at (13,0)
    {\iflangja{m の待ちキュー}{wait queue of mutex m}};
  \node[st, minimum width=28mm, text width=26mm, fill=white, thick] (cs) at (58,0)
    {\iflangja{クリティカルセクション\\（m を保持）}{critical section\\(owns m)}};
  \node[st, minimum width=26mm, text width=24mm] (cq) at (103,0)
    {\iflangja{c の待ちキュー}{wait queue of condition c}};
  \node[font=\footnotesize\ttfamily, anchor=south] (lk) at (35,21) {Lock(m)};
  \draw (lk.south) -- (35,16);
  \draw[->] (35,16) -| node[lab, pos=0.4, above] {\iflangja{m は使用中}{m held}} (mq.north);
  \draw[->] (35,16) -| node[lab, pos=0.5, above] {\iflangja{m は空き}{m free}} (cs.north);
  \draw[->] (mq) -- node[lab, above] {\iflangja{m を獲得}{acquire m}} (cs);
  \draw[->] (cs) -- node[lab, above] {\texttt{Wait(m,c)}}
                    node[lab, below] {\iflangja{m を解放}{releases m}} (cq);
  \draw[->] ([xshift=8mm]cs.north) |- node[lab, pos=0.75, above] {\texttt{Unlock(m)}} (94,16);
  % signalled: re-acquire m directly if it is free, otherwise queue on m
  \draw (cq.south) -- node[lab, left, pos=0.55] {\texttt{Signal(c)} / \texttt{Broadcast(c)}} (103,-15);
  \draw[->] (103,-15) -| node[lab, pos=0.75, left] {\iflangja{m は空き}{m free}} ([xshift=8mm]cs.south);
  \draw[->] (103,-15) -- (103,-21) -| node[lab, pos=0.25, below] {\iflangja{m は使用中}{m held}} (mq.south);
\end{tikzpicture}

  \caption{条件 \texttt{c} で待つスレッドはミューテックス \texttt{m} を
    解放する．signal されると，スレッドはまず \texttt{m} を再獲得しなければ
    ならない．\texttt{m} が保持されていれば，ミューテックスの待ちキューに
    移る．}
  \label{fig:l03-cv-queues}
\end{figure}

消費者が \texttt{if} ではなく \texttt{while} ループでリストを検査するのは
このためである．\texttt{Signal} から消費者がミューテックスを再獲得する
までの間に，他のスレッドが先にミューテックスを獲得してリストを変更する
かもしれない---例えば，リストを空にする2番目の消費者である．条件変数への
signal は，条件が成り立っているかもしれないことを伝えるだけである．
起こされたスレッドがすぐにミューテックスを獲得できる保証はなく，獲得
したら条件を再検査しなければならない．

\needspace{7\baselineskip}
\begin{keypoint}
  \texttt{Wait} は必ずミューテックスを保持した状態で，条件を再検査する
  \texttt{while} ループの中で呼び出す．signal または broadcast は，条件
  が依存する状態を変更した後に行う．
\end{keypoint}

\section{リーダ・ライタ問題}
\label{sec:l03-rw}

\term[りーだらいたもんだい]{リーダ・ライタ問題}[readers-writers problem]%
では，一部のスレッド（リーダ）はファイルなどの共有資源を読むだけで，他の
スレッド（ライタ）はそれを変更する．\source{P2L2，リーダ/ライタ問題と
その例，（代理変数を用いた）クリティカルセクションの構造；OSTEP
ch.~31．}任意の数のリーダが同時に資源にアクセスしてよいが，ライタは排他的
なアクセスを必要とする．すなわち，ライタは高々1つで，書き込み中はリーダも
いてはならない．

資源を1つのミューテックスで保護するのは正しいが，制限が強すぎる．リーダ
同士も互いに排除してしまうからである．アクセス規則は，活動中のリーダと
ライタの数に依存する．
\begin{itemize}
  \item 活動中のリーダもライタもいない: リーダもライタも入ってよい．
  \item 活動中のリーダが1つ以上いる: 別のリーダは入ってよいが，ライタは
    入れない．
  \item 活動中のライタがいる: リーダもライタも入れない．
\end{itemize}
これらの規則は1つのカウンタで表現できる．資源は，
\texttt{resource\_counter} $=0$ のとき\emph{空き}，
\texttt{resource\_counter} $=n>0$ のとき $n$ 個のリーダによって
\emph{読み出し中}，\texttt{resource\_counter} $=-1$ のとき
\emph{書き込み中}である（\cref{fig:l03-rw-states}）．

\begin{figure}[htbp]
  \centering
  % States of the shared resource in the readers-writers solution.
% Usage: % States of the shared resource in the readers-writers solution.
% Usage: % States of the shared resource in the readers-writers solution.
% Usage: \input{figures/l03-rw-states}   (width ~116 mm)
\begin{tikzpicture}[x=1mm, y=1mm, font=\footnotesize\sffamily,
  st/.style={draw, rounded corners=5pt, fill=ln-box, align=center,
             minimum width=26mm, minimum height=11mm, inner sep=2pt},
  >={Stealth[length=1.8mm]}, lab/.style={font=\scriptsize\sffamily, align=center}]
  \node[st] (wr) at (13,0)  {\iflangja{書き込み中}{writing}\\ \texttt{counter} $=-1$};
  \node[st, fill=white] (fr) at (58,0)  {\iflangja{空き}{free}\\ \texttt{counter} $=0$};
  \node[st] (rd) at (103,0) {\iflangja{読み出し中}{reading}\\ \texttt{counter} $=n>0$};
  \draw[->] (fr.north west) to[bend right=30]
    node[lab, above] {\iflangja{ライタが入る}{writer enters}} (wr.north east);
  \draw[->] (wr.south east) to[bend right=30]
    node[lab, below] {\iflangja{ライタが出る}{writer exits}} (fr.south west);
  \draw[->] (fr.north east) to[bend left=30]
    node[lab, above] {\iflangja{リーダが入る}{reader enters}} (rd.north west);
  \draw[->] (rd.south west) to[bend left=30]
    node[lab, below] {\iflangja{最後のリーダが出る}{last reader exits}} (fr.south east);
  \path[->] (rd) edge[loop above, min distance=7mm, in=60, out=120]
    node[lab, above] {\iflangja{他のリーダが入る/出る}{another reader enters or exits}} (rd);
\end{tikzpicture}
   (width ~116 mm)
\begin{tikzpicture}[x=1mm, y=1mm, font=\footnotesize\sffamily,
  st/.style={draw, rounded corners=5pt, fill=ln-box, align=center,
             minimum width=26mm, minimum height=11mm, inner sep=2pt},
  >={Stealth[length=1.8mm]}, lab/.style={font=\scriptsize\sffamily, align=center}]
  \node[st] (wr) at (13,0)  {\iflangja{書き込み中}{writing}\\ \texttt{counter} $=-1$};
  \node[st, fill=white] (fr) at (58,0)  {\iflangja{空き}{free}\\ \texttt{counter} $=0$};
  \node[st] (rd) at (103,0) {\iflangja{読み出し中}{reading}\\ \texttt{counter} $=n>0$};
  \draw[->] (fr.north west) to[bend right=30]
    node[lab, above] {\iflangja{ライタが入る}{writer enters}} (wr.north east);
  \draw[->] (wr.south east) to[bend right=30]
    node[lab, below] {\iflangja{ライタが出る}{writer exits}} (fr.south west);
  \draw[->] (fr.north east) to[bend left=30]
    node[lab, above] {\iflangja{リーダが入る}{reader enters}} (rd.north west);
  \draw[->] (rd.south west) to[bend left=30]
    node[lab, below] {\iflangja{最後のリーダが出る}{last reader exits}} (fr.south east);
  \path[->] (rd) edge[loop above, min distance=7mm, in=60, out=120]
    node[lab, above] {\iflangja{他のリーダが入る/出る}{another reader enters or exits}} (rd);
\end{tikzpicture}
   (width ~116 mm)
\begin{tikzpicture}[x=1mm, y=1mm, font=\footnotesize\sffamily,
  st/.style={draw, rounded corners=5pt, fill=ln-box, align=center,
             minimum width=26mm, minimum height=11mm, inner sep=2pt},
  >={Stealth[length=1.8mm]}, lab/.style={font=\scriptsize\sffamily, align=center}]
  \node[st] (wr) at (13,0)  {\iflangja{書き込み中}{writing}\\ \texttt{counter} $=-1$};
  \node[st, fill=white] (fr) at (58,0)  {\iflangja{空き}{free}\\ \texttt{counter} $=0$};
  \node[st] (rd) at (103,0) {\iflangja{読み出し中}{reading}\\ \texttt{counter} $=n>0$};
  \draw[->] (fr.north west) to[bend right=30]
    node[lab, above] {\iflangja{ライタが入る}{writer enters}} (wr.north east);
  \draw[->] (wr.south east) to[bend right=30]
    node[lab, below] {\iflangja{ライタが出る}{writer exits}} (fr.south west);
  \draw[->] (fr.north east) to[bend left=30]
    node[lab, above] {\iflangja{リーダが入る}{reader enters}} (rd.north west);
  \draw[->] (rd.south west) to[bend left=30]
    node[lab, below] {\iflangja{最後のリーダが出る}{last reader exits}} (fr.south east);
  \path[->] (rd) edge[loop above, min distance=7mm, in=60, out=120]
    node[lab, above] {\iflangja{他のリーダが入る/出る}{another reader enters or exits}} (rd);
\end{tikzpicture}

  \caption{共有資源の状態と，それに対応する代理変数の値．\texttt{counter}
    は \texttt{resource\_counter} の略である．}
  \label{fig:l03-rw-states}
\end{figure}

このカウンタは\term[だいりへんすう]{代理変数}[proxy variable]の一例で
ある．代理変数とは，ミューテックスで保護され，資源の状態を表す共有変数で
ある．スレッドがミューテックスを保持するのは代理変数を検査・更新する間だけ
であり，実際の読み書きは，代理変数が強制する方針のもとでミューテックスの
外で行われる．解は1つのミューテックスと2つの条件変数を用いる．
\begin{lstlisting}[language=C, numbers=none]
Mutex counter_mutex;
Condition read_phase, write_phase;
int resource_counter = 0;

// リーダ
Lock(counter_mutex) {
  while (resource_counter == -1)
    Wait(counter_mutex, read_phase);
  resource_counter++;
} // unlock
// ... データを読む ...
Lock(counter_mutex) {
  resource_counter--;
  if (resource_counter == 0)
    Signal(write_phase);
} // unlock
\end{lstlisting}
\needspace{14\baselineskip}
\begin{lstlisting}[language=C, numbers=none]
// ライタ
Lock(counter_mutex) {
  while (resource_counter != 0)
    Wait(counter_mutex, write_phase);
  resource_counter = -1;
} // unlock
// ... データを書く ...
Lock(counter_mutex) {
  resource_counter = 0;
  Broadcast(read_phase);
  Signal(write_phase);
} // unlock
\end{lstlisting}

退出処理でどの通知を用いるかは，誰が待っている可能性があるかから決まる．
\begin{itemize}
  \item 他のリーダが残っている間に出ていく\textbf{リーダ}は，待っている
    スレッドに関係する状態を何も変えない．最後のリーダが出ていくとき，
    待っている可能性があるのはライタだけであり（リーダが待つのはライタが
    活動中のときだけである），入れるのはそのうち1つだけなので，
    \texttt{write\_phase} への \texttt{Signal} で十分である．
  \item 出ていく\textbf{ライタ}は，どちらの種類のスレッドに対しても資源を
    空ける．待っているリーダはすべて一緒に入ってよいので
    \texttt{Broadcast(read\_phase)} を，ライタは1つしか入れないので
    \texttt{Signal(write\_phase)} を行う．実際にどのスレッドが先に入るか
    は，どれが先に \texttt{counter\_mutex} を獲得するかに依存する．リーダ
    が先に入れば，起こされた他のリーダも入り，ライタは \texttt{while}
    ループでカウンタ $\neq 0$ を見て再び待つ．ライタが先に入れば，リーダは
    カウンタ $=-1$ を見て再び待つ．
\end{itemize}

\begin{example}
  \cref{tab:l03-rw-trace} は，2つのリーダ，それらの読み出し中に到着する
  ライタ，ライタの書き込み中に到着する3番目のリーダをトレースしたもので
  ある．
\end{example}

\begin{table}[htbp]
  \centering
  \caption{リーダ・ライタ問題の解のトレース．各ステップの後の代理変数の
    値を示す．}
  \label{tab:l03-rw-trace}
  \small
  \begin{tabular}{@{}clr@{}}
    \toprule
    ステップ & イベント & カウンタ \\
    \midrule
    1 & R1 が入る                                               & $1$ \\
    2 & R2 が入る                                               & $2$ \\
    3 & W1 はカウンタ $\neq 0$ を見て \texttt{write\_phase} で待つ & $2$ \\
    4 & R1 が出る．カウンタ $\neq 0$ なので通知しない             & $1$ \\
    5 & R2 が出る．カウンタ $=0$ なので \texttt{write\_phase} に signal & $0$ \\
    6 & W1 がミューテックスを再獲得し，再検査して入る             & $-1$ \\
    7 & R3 はカウンタ $=-1$ を見て \texttt{read\_phase} で待つ    & $-1$ \\
    8 & W1 が出る．\texttt{read\_phase} に broadcast，\texttt{write\_phase} に signal & $0$ \\
    9 & R3 がミューテックスを再獲得し，再検査して入る             & $1$ \\
    \bottomrule
  \end{tabular}
\end{table}

\needspace{14\baselineskip}
\subsection{クリティカルセクションの構造}

リーダ・ライタ問題の解は，ミューテックスと条件変数を組み合わせて使うコード
の一般的な形を示している．待つ必要がありうるクリティカルセクションは次の
ように書く．
\begin{lstlisting}[language=C, numbers=none]
Lock(mutex) {
  while (!predicate_indicating_access_ok)
    Wait(mutex, cond_var);
  update state => update predicate;
  Signal and/or Broadcast(
      cond_var_with_correct_waiting_threads);
} // unlock
\end{lstlisting}
\needspace{9\baselineskip}
代理変数を用いる場合は，資源に対する操作をミューテックスの外に出し，全体を
3つの部分に分ける．
\begin{lstlisting}[language=C, numbers=none]
// クリティカルセクションに入る
Lock(mutex) {
  while (!predicate_for_access)
    Wait(mutex, cond_var);
  update predicate;
} // unlock

perform critical operation  // 例: 共有ファイルの読み書き

// クリティカルセクションから出る
Lock(mutex) {
  update predicate;
  Signal/Broadcast(cond_var);
} // unlock
\end{lstlisting}

\needspace{6\baselineskip}
\begin{keypoint}
  ミューテックスだけで強制できる方針は「一度に1つのスレッド」の1つしか
  ない．短いクリティカルセクションで更新される代理変数を用いると，
  「任意の数のリーダ，または1つのライタ」のようなより複雑なアクセス方針を，
  同じミューテックスと条件変数のプリミティブから構築できる．
\end{keypoint}

\section{よくある誤りとスプリアスウェイクアップ}
\label{sec:l03-pitfalls}

ミューテックスと条件変数を使うプログラムでは，いくつかの誤りが繰り返し
現れる．\source{P2L2，よくある誤り，スプリアスウェイクアップ；Birrell
\cite{birrell1989}．}
\begin{itemize}
  \item どのミューテックスと条件変数がどの資源に属するかを見失う．変数を
    宣言する箇所に，例えば \texttt{Mutex m1; // file1 を保護} のように
    対応関係を記録すれば，これを避けられる．
  \item ロックやアンロックを忘れる，あるいは誤ったミューテックスをロック・
    アンロックする．コンパイラや解析ツールの中には，このような場合に警告
    を出せるものがある．
  \item 同じ資源に複数のミューテックスを使う．ファイルのリーダが
    \texttt{m1} を，ライタが \texttt{m2} をロックすると，読み出しと
    書き込みが並行に実行されてしまう．
  \item 誤った条件変数に signal する．
  \item \texttt{Broadcast} が必要なところで \texttt{Signal} を使う．
    起こされるのは1つのスレッドだけで，他のスレッドはいつまでも待ち続ける
    かもしれない．
  \item 通知によって実行順序が決まると仮定する．signal を発行した順序は，
    起こされたスレッドが実行される順序を決めない．優先度を守らなければ
    ならない場合は，別の機構が必要である．
\end{itemize}

さらに2つの問題を詳しく見る．性能を損なうスプリアスウェイクアップと，
プログラムを停止させるデッドロック（\cref{sec:l03-deadlock}）である．

\subsection{スプリアスウェイクアップ}

通知の中には，それを利用できないスレッドを起こしてしまうものがある．

\begin{definition}[スプリアスウェイクアップ]
  \term[すぷりあすうぇいくあっぷ]{スプリアスウェイクアップ}[spurious wake-up]%
  とは，実行されても先に進めないスレッドを起こすことであり，それによって
  生じるコンテキストスイッチは無駄になる．
\end{definition}

よくある原因は，ミューテックスを保持したまま signal することである．
\cref{sec:l03-rw}のライタの退出処理では，broadcast はクリティカル
セクションの中で発行される．起こされたリーダは \texttt{read\_phase} の
待ちキューから出るが，ライタがまだ保持している \texttt{counter\_mutex} を
再獲得するまで \texttt{Wait} から戻れない．それらは単にミューテックスの
待ちキューに移されるだけである（\cref{fig:l03-cv-queues}）．
\margin{POSIX では，signal がまったくないのに wait から戻ることも
スプリアスウェイクアップと呼ぶ．\texttt{while} ループはどちらにも対処
する．}

\needspace{9\baselineskip}
通知が共有状態に依存しない場合，対策は通知をアンロックの後に移すことで
ある．
\begin{lstlisting}[language=C, numbers=none]
// ライタの退出処理（通知はミューテックスの外）
Lock(counter_mutex) {
  resource_counter = 0;
} // unlock
Broadcast(read_phase);
Signal(write_phase);
\end{lstlisting}

これは常に可能なわけではない．リーダの退出処理では，signal するかどうか
の判断が \texttt{resource\_counter} に依存し，この変数はミューテックスを
保持している間しか読めないので，\texttt{Signal} を単純にアンロックの後に
移すことはできない．移せるのは，クリティカルセクションの中で判断を局所
変数に記録し，その後で signal する場合だけである．スプリアスウェイクアップ
が正しさに影響することはなく（先に進めないスレッドは \texttt{while}
ループによって待ちに戻される），影響するのは性能だけである．

\section{デッドロック}
\label{sec:l03-deadlock}

一部のミューテックスを保持したまま別のミューテックスを要求するスレッドは，
互いを永久にブロックしうる．\source{P2L2，デッドロック；OSTEP ch.~32；
Silberschatz ら ch.~8．}

\begin{definition}[デッドロック]
  \term[でっどろっく]{デッドロック}[deadlock]とは，競合する2つ以上の
  スレッドがそれぞれ他のスレッドの完了を待っているが，どれも決して完了
  しない状況である．
\end{definition}

\needspace{10\baselineskip}
最も単純な場合は，2つのスレッドと2つのミューテックスが関わるものである．
どちらのスレッドも \texttt{foo(A, B)} を実行するために資源 A と B を必要と
するが，それらを保護するミューテックスを逆の順序でロックする．
\begin{lstlisting}[language=C, numbers=none]
// スレッド T1             // スレッド T2
Lock(m_A) {               Lock(m_B) {
  Lock(m_B) {               Lock(m_A) {
    foo(A, B);                foo(A, B);
  }                         }
}                         }
\end{lstlisting}
T1 が \texttt{m\_A} を，T2 が \texttt{m\_B} を，どちらも2番目の
\texttt{Lock} に到達する前に獲得すると，T1 は \texttt{m\_B} を待って
ブロックし，T2 は \texttt{m\_A} を待ってブロックする．どちらも先に進めず，
どちらも相手が必要とするミューテックスを決して解放しない．

この状況は\term[まちぐらふ]{待ちグラフ}[wait graph]で記述される．待ち
グラフは，資源を待っているスレッドからその資源を所有するスレッドへの辺を
持つ．\cref{fig:l03-wait-graph} は，ミューテックスを明示的に描いてこの例を
示したものである．デッドロックは待ちグラフの閉路に対応する．所有者が1つ
だけであるロックについては，閉路はデッドロックの必要十分条件である．

\begin{figure}[htbp]
  \centering
  % Wait graph of the two-thread, two-lock deadlock.
% Usage: % Wait graph of the two-thread, two-lock deadlock.
% Usage: % Wait graph of the two-thread, two-lock deadlock.
% Usage: \input{figures/l03-wait-graph}   (width ~80 mm)
\begin{tikzpicture}[x=1mm, y=1mm, font=\footnotesize\sffamily,
  th/.style={draw, circle, fill=ln-box, minimum size=9mm, inner sep=0pt},
  lk/.style={draw, rectangle, fill=white, thick, minimum size=7mm, inner sep=0pt},
  >={Stealth[length=1.8mm]}, lab/.style={font=\scriptsize\sffamily},
  holds/.style={->, thick}, waits/.style={->, thick, dashed, ln-caution}]
  \node[th] (t1) at (0,0)   {T1};
  \node[th] (t2) at (60,0)  {T2};
  \node[lk] (a)  at (30,14) {\texttt{m\_A}};
  \node[lk] (b)  at (30,-14) {\texttt{m\_B}};
  \draw[holds] (a) -- node[lab, above left]  {\iflangja{T1 が保持}{held by T1}} (t1);
  \draw[waits] (t1) -- node[lab, below left] {\iflangja{T1 が待つ}{T1 waits for}} (b);
  \draw[holds] (b) -- node[lab, below right] {\iflangja{T2 が保持}{held by T2}} (t2);
  \draw[waits] (t2) -- node[lab, above right] {\iflangja{T2 が待つ}{T2 waits for}} (a);
\end{tikzpicture}
   (width ~80 mm)
\begin{tikzpicture}[x=1mm, y=1mm, font=\footnotesize\sffamily,
  th/.style={draw, circle, fill=ln-box, minimum size=9mm, inner sep=0pt},
  lk/.style={draw, rectangle, fill=white, thick, minimum size=7mm, inner sep=0pt},
  >={Stealth[length=1.8mm]}, lab/.style={font=\scriptsize\sffamily},
  holds/.style={->, thick}, waits/.style={->, thick, dashed, ln-caution}]
  \node[th] (t1) at (0,0)   {T1};
  \node[th] (t2) at (60,0)  {T2};
  \node[lk] (a)  at (30,14) {\texttt{m\_A}};
  \node[lk] (b)  at (30,-14) {\texttt{m\_B}};
  \draw[holds] (a) -- node[lab, above left]  {\iflangja{T1 が保持}{held by T1}} (t1);
  \draw[waits] (t1) -- node[lab, below left] {\iflangja{T1 が待つ}{T1 waits for}} (b);
  \draw[holds] (b) -- node[lab, below right] {\iflangja{T2 が保持}{held by T2}} (t2);
  \draw[waits] (t2) -- node[lab, above right] {\iflangja{T2 が待つ}{T2 waits for}} (a);
\end{tikzpicture}
   (width ~80 mm)
\begin{tikzpicture}[x=1mm, y=1mm, font=\footnotesize\sffamily,
  th/.style={draw, circle, fill=ln-box, minimum size=9mm, inner sep=0pt},
  lk/.style={draw, rectangle, fill=white, thick, minimum size=7mm, inner sep=0pt},
  >={Stealth[length=1.8mm]}, lab/.style={font=\scriptsize\sffamily},
  holds/.style={->, thick}, waits/.style={->, thick, dashed, ln-caution}]
  \node[th] (t1) at (0,0)   {T1};
  \node[th] (t2) at (60,0)  {T2};
  \node[lk] (a)  at (30,14) {\texttt{m\_A}};
  \node[lk] (b)  at (30,-14) {\texttt{m\_B}};
  \draw[holds] (a) -- node[lab, above left]  {\iflangja{T1 が保持}{held by T1}} (t1);
  \draw[waits] (t1) -- node[lab, below left] {\iflangja{T1 が待つ}{T1 waits for}} (b);
  \draw[holds] (b) -- node[lab, below right] {\iflangja{T2 が保持}{held by T2}} (t2);
  \draw[waits] (t2) -- node[lab, above right] {\iflangja{T2 が待つ}{T2 waits for}} (a);
\end{tikzpicture}

  \caption{このデッドロックの待ちグラフ．T1 から辺をたどると，
    \texttt{m\_B}，T2，\texttt{m\_A} を経て T1 に戻る．}
  \label{fig:l03-wait-graph}
\end{figure}

\subsection{例のデッドロックの回避}

例からデッドロックを取り除く方法は3つある．
\begin{enumerate}
  \item \textbf{B をロックする前に A をアンロックする．}一度に1つの
    ミューテックスだけを保持すれば（細粒度のロック）閉路は生じえないが，
    ここでは使えない．スレッドは両方の資源を同時に必要とするからである．
  \item \textbf{すべてのロックを最初に獲得する．}開始前に必要なすべての
    ロックを一度に（アトミックに，すなわちすべてを獲得するか1つも獲得
    しないかのどちらかで，1つでも得られなければ保持しているものを解放
    して）獲得し，最後にすべて解放するスレッドはデッドロックしえない．その極端な場合が，すべての資源に対する1つの
    「メガロック」である．並列性を制限するので，多くのアプリケーション
    にとってこれは制限が強すぎる．
  \item \textbf{ロックの順序を守る．}すべてのスレッドがミューテックスを
    同じ順序，すなわち先に \texttt{m\_A}，次に \texttt{m\_B} の順で獲得
    すれば，デッドロックは起こりえない．
\end{enumerate}

3番目の方法である\term[ろっくじゅんじょ]{ロック順序}[lock ordering]は，
一般に通用する．ミューテックスに番号を付け，すべてのスレッドに番号の昇順
で獲得させる．スレッド $T_1,\dots,T_k$ が閉路を形成し，各 $T_j$ が次の
スレッド $T_{j+1}$（ただし $T_{k+1}=T_1$）の所有するミューテックス $M_j$
を待っていると仮定する．スレッド $T_{j+1}$ は $M_j$ を保持して $M_{j+1}$
を要求しているので，$M_{j+1}>M_j$ である．すると番号は閉路を一周する間
ずっと狭義に増加しながら $M_1$ に戻ることになり，これは不可能である．
難しさは実践面にある．多数のミューテックスを持つ大規模なプログラムでは，
すべてのコード経路が順序を守ることを保証するのは難しく，ロックの獲得順序
を検査するツールが役立つ．

\subsection{デッドロックへの一般的な対処}

オペレーティングシステムとアプリケーションがデッドロックに対処する方法は
3つある．
\begin{description}
  \item[防止] \term[でっどろっくぼうし]{デッドロック防止}[deadlock prevention]%
    では，閉路の形成を決して許さない．コードを変更して閉路を不可能にする
    （スレッドが別のロックを要求する前に保持しているロックを解放する，
    あるいはすべてのスレッドがロック順序に従う）か，すべてのロック要求を
    実行時に検査し，それを許可すると後に閉路につながりうる場合には要求を
    遅らせる．実行時の検査はコストが高く，各スレッドがどのロックを要求
    しうるかを事前に知る必要がある．教科書ではこの変種を
    \emph{デッドロック回避}と呼び，\emph{防止}という語を構造的な規則に
    限って用いる．
  \item[検出と回復]
    \term[でっどろっくのけんしゅつとかいふく]{デッドロックの検出と回復}[deadlock detection and recovery]%
    では，システムが定期的に待ちグラフを解析して閉路を探し，見つかった
    場合には関係するスレッドの一部の実行をロールバックする．ロール
    バックには以前の時点を復元するのに十分な保存状態が必要であり，
    スレッドがすでに外部からの入力を消費したり外部への出力を生成したり
    している場合には不可能である．
  \item[何もしない] \term[だちょうあるごりずむ]{ダチョウアルゴリズム}[ostrich algorithm]%
    は，デッドロックは稀であり他の方法はコストが高いという理由で，
    デッドロックを完全に無視する．実際にデッドロックが起きたら，システム
    を再起動する．
\end{description}

\needspace{7\baselineskip}
\begin{keypoint}
  デッドロックは待ちグラフの閉路である．すべてのミューテックスを1つの
  大域的な順序で獲得すれば，実行時コストなしにそのような閉路を防止できる．
  すべてのロック要求の実行時検査やロールバックによる回復は，汎用的だが
  コストが高い．
\end{keypoint}

\section{カーネルレベルスレッドとユーザレベルスレッド}
\label{sec:l03-klt-ult}

スレッドはシステムの2つのレベルに存在する．\source{P2L2，カーネルレベル
スレッドとユーザレベルスレッド，マルチスレッドモデル，マルチスレッドの
スコープ；Silberschatz ら ch.~4--5．}オペレーティングシステムのカーネル
自体がマルチスレッドでありうるし，アプリケーションもプロセス内に
スレッドを生成できる．問題は，この2種類のスレッドが互いにどう関係するか
である．
\begin{definition}[カーネルレベルスレッドとユーザレベルスレッド]
  \term[かーねるれべるすれっど]{カーネルレベルスレッド}[kernel-level thread]%
  とは，カーネルから見え，CPU スケジューラなどのカーネルの構成要素に
  よって管理されるスレッドである．
  \term[ゆーざれべるすれっど]{ユーザレベルスレッド}[user-level thread]%
  とは，ユーザレベルのスレッドライブラリを通じて生成される，アプリケー
  ションから見たスレッドである．マルチスレッドモデルに応じて，ライブラリ
  はそれを自ら管理するか，操作をカーネルに渡す．
\end{definition}

カーネルレベルスレッドの中には，アプリケーションのユーザレベルスレッドを
支えるために存在するものもあれば，デーモンなどのオペレーティングシステム
サービスを実行するものもある．CPU スケジューラは，どのカーネルレベル
スレッドをどの CPU で実行するかを決める．一方ユーザレベルスレッドは，
カーネルレベルスレッドに対応付けられ，スケジューラがそのカーネルレベル
スレッドを CPU 上で実行するときにのみ実行できる．

\subsection{マルチスレッドモデル}

ユーザレベルスレッドをカーネルレベルスレッドにどう対応付けるかは，3つの
マルチスレッドモデルで記述される（\cref{fig:l03-mt-models}）．

\begin{figure}[htbp]
  \centering
  % Multithreading models: mapping of user-level threads to kernel-level threads.
% Usage: % Multithreading models: mapping of user-level threads to kernel-level threads.
% Usage: % Multithreading models: mapping of user-level threads to kernel-level threads.
% Usage: \input{figures/l03-mt-models}   (width ~118 mm)
\begin{tikzpicture}[x=0.96mm, y=1mm, font=\footnotesize\sffamily,
  ut/.style={draw, circle, fill=white, minimum size=4mm, inner sep=0pt},
  kt/.style={draw, rectangle, fill=ln-box, minimum size=4mm, inner sep=0pt},
  lib/.style={draw, rounded corners=1pt, fill=ln-box, minimum height=4mm,
              inner sep=1pt, font=\scriptsize\sffamily},
  ttl/.style={font=\footnotesize\sffamily\bfseries},
  lab/.style={font=\scriptsize\sffamily, text=ln-muted}]
  \draw[dashed, ln-rule] (0,10) -- (118,10);
  \node[lab, anchor=south west] at (0,10.5) {\iflangja{ユーザ}{user}};
  \node[lab, anchor=north west] at (0,9.5)  {\iflangja{カーネル}{kernel}};
  % ---- one-to-one
  \node[ttl] at (27,32) {\iflangja{一対一}{one-to-one}};
  \foreach \x in {17,27,37} {
    \node[ut] (u) at (\x,24) {};
    \node[kt] (k) at (\x,3) {};
    \draw (u) -- (k);
  }
  % ---- many-to-one
  \node[ttl] at (65,32) {\iflangja{多対一}{many-to-one}};
  \node[lib, minimum width=28mm] (l2) at (65,16) {\iflangja{スレッドライブラリ}{thread library}};
  \foreach \x in {55,65,75} { \node[ut] (u) at (\x,24) {}; \draw (u) -- (\x,18); }
  \node[kt] (k2) at (65,3) {};
  \draw (l2) -- (k2);
  % ---- many-to-many
  \node[ttl] at (103,32) {\iflangja{多対多}{many-to-many}};
  \node[lib, minimum width=24mm] (l3) at (99,16) {\iflangja{ライブラリ}{library}};
  \foreach \x in {90,99,108} { \node[ut] (u) at (\x,24) {}; \draw (u) -- (\x,18); }
  \node[kt] (k3a) at (93,3) {};
  \node[kt] (k3b) at (105,3) {};
  \draw (l3.south) ++(-6,0) -- (k3a);
  \draw (l3.south) ++(6,0) -- (k3b);
  \node[ut, thick, draw=ln-accent] (ub) at (116,24) {};
  \node[kt, thick, draw=ln-accent] (kb) at (116,3) {};
  \draw[very thick, ln-accent] (ub) -- (kb);
  \node[lab, anchor=north, text=ln-accent] at (116,0.5) {\iflangja{バウンド}{bound}};
\end{tikzpicture}
   (width ~118 mm)
\begin{tikzpicture}[x=0.96mm, y=1mm, font=\footnotesize\sffamily,
  ut/.style={draw, circle, fill=white, minimum size=4mm, inner sep=0pt},
  kt/.style={draw, rectangle, fill=ln-box, minimum size=4mm, inner sep=0pt},
  lib/.style={draw, rounded corners=1pt, fill=ln-box, minimum height=4mm,
              inner sep=1pt, font=\scriptsize\sffamily},
  ttl/.style={font=\footnotesize\sffamily\bfseries},
  lab/.style={font=\scriptsize\sffamily, text=ln-muted}]
  \draw[dashed, ln-rule] (0,10) -- (118,10);
  \node[lab, anchor=south west] at (0,10.5) {\iflangja{ユーザ}{user}};
  \node[lab, anchor=north west] at (0,9.5)  {\iflangja{カーネル}{kernel}};
  % ---- one-to-one
  \node[ttl] at (27,32) {\iflangja{一対一}{one-to-one}};
  \foreach \x in {17,27,37} {
    \node[ut] (u) at (\x,24) {};
    \node[kt] (k) at (\x,3) {};
    \draw (u) -- (k);
  }
  % ---- many-to-one
  \node[ttl] at (65,32) {\iflangja{多対一}{many-to-one}};
  \node[lib, minimum width=28mm] (l2) at (65,16) {\iflangja{スレッドライブラリ}{thread library}};
  \foreach \x in {55,65,75} { \node[ut] (u) at (\x,24) {}; \draw (u) -- (\x,18); }
  \node[kt] (k2) at (65,3) {};
  \draw (l2) -- (k2);
  % ---- many-to-many
  \node[ttl] at (103,32) {\iflangja{多対多}{many-to-many}};
  \node[lib, minimum width=24mm] (l3) at (99,16) {\iflangja{ライブラリ}{library}};
  \foreach \x in {90,99,108} { \node[ut] (u) at (\x,24) {}; \draw (u) -- (\x,18); }
  \node[kt] (k3a) at (93,3) {};
  \node[kt] (k3b) at (105,3) {};
  \draw (l3.south) ++(-6,0) -- (k3a);
  \draw (l3.south) ++(6,0) -- (k3b);
  \node[ut, thick, draw=ln-accent] (ub) at (116,24) {};
  \node[kt, thick, draw=ln-accent] (kb) at (116,3) {};
  \draw[very thick, ln-accent] (ub) -- (kb);
  \node[lab, anchor=north, text=ln-accent] at (116,0.5) {\iflangja{バウンド}{bound}};
\end{tikzpicture}
   (width ~118 mm)
\begin{tikzpicture}[x=0.96mm, y=1mm, font=\footnotesize\sffamily,
  ut/.style={draw, circle, fill=white, minimum size=4mm, inner sep=0pt},
  kt/.style={draw, rectangle, fill=ln-box, minimum size=4mm, inner sep=0pt},
  lib/.style={draw, rounded corners=1pt, fill=ln-box, minimum height=4mm,
              inner sep=1pt, font=\scriptsize\sffamily},
  ttl/.style={font=\footnotesize\sffamily\bfseries},
  lab/.style={font=\scriptsize\sffamily, text=ln-muted}]
  \draw[dashed, ln-rule] (0,10) -- (118,10);
  \node[lab, anchor=south west] at (0,10.5) {\iflangja{ユーザ}{user}};
  \node[lab, anchor=north west] at (0,9.5)  {\iflangja{カーネル}{kernel}};
  % ---- one-to-one
  \node[ttl] at (27,32) {\iflangja{一対一}{one-to-one}};
  \foreach \x in {17,27,37} {
    \node[ut] (u) at (\x,24) {};
    \node[kt] (k) at (\x,3) {};
    \draw (u) -- (k);
  }
  % ---- many-to-one
  \node[ttl] at (65,32) {\iflangja{多対一}{many-to-one}};
  \node[lib, minimum width=28mm] (l2) at (65,16) {\iflangja{スレッドライブラリ}{thread library}};
  \foreach \x in {55,65,75} { \node[ut] (u) at (\x,24) {}; \draw (u) -- (\x,18); }
  \node[kt] (k2) at (65,3) {};
  \draw (l2) -- (k2);
  % ---- many-to-many
  \node[ttl] at (103,32) {\iflangja{多対多}{many-to-many}};
  \node[lib, minimum width=24mm] (l3) at (99,16) {\iflangja{ライブラリ}{library}};
  \foreach \x in {90,99,108} { \node[ut] (u) at (\x,24) {}; \draw (u) -- (\x,18); }
  \node[kt] (k3a) at (93,3) {};
  \node[kt] (k3b) at (105,3) {};
  \draw (l3.south) ++(-6,0) -- (k3a);
  \draw (l3.south) ++(6,0) -- (k3b);
  \node[ut, thick, draw=ln-accent] (ub) at (116,24) {};
  \node[kt, thick, draw=ln-accent] (kb) at (116,3) {};
  \draw[very thick, ln-accent] (ub) -- (kb);
  \node[lab, anchor=north, text=ln-accent] at (116,0.5) {\iflangja{バウンド}{bound}};
\end{tikzpicture}

  \caption{マルチスレッドモデル．円はユーザレベルスレッド，四角はカーネル
    レベルスレッドである．多対多モデルでは，バウンドスレッドは専用の
    カーネルレベルスレッドを持ち続ける．}
  \label{fig:l03-mt-models}
\end{figure}

\begin{description}
  \item[一対一] \term[いったいいちもでる]{一対一モデル}[one-to-one model]%
    では，各ユーザレベルスレッドが専用のカーネルレベルスレッドを持つ．
    カーネルはスレッド，その同期，ブロッキングを把握・理解し，それらを
    複数の CPU にスケジュールできる．一方で，スレッド操作はすべて
    システムコールでカーネルを経由するためコストが高い．アプリケーション
    はカーネルがサポートする方針とスレッド数の上限に制約され，この
    サポートを提供するカーネルにしか移植できない．
  \item[多対一] \term[たたいいちもでる]{多対一モデル}[many-to-one model]%
    では，プロセスのすべてのユーザレベルスレッドが1つのカーネルレベル
    スレッドに対応付けられ，ユーザレベルのスレッドライブラリが，その上で
    どのユーザレベルスレッドを実行するかを随時決める．スレッド管理に
    システムコールは不要で，カーネルの制限や方針にも依存しないので，完全
    に移植可能である．しかしカーネルはアプリケーションが何を必要として
    いるかを知る手段を持たず，1つのユーザレベルスレッドが I/O で
    ブロックするとカーネルレベルスレッドがブロックし，それとともに
    プロセス全体がブロックする．
  \item[多対多] \term[たたいたもでる]{多対多モデル}[many-to-many model]%
    では，一部のユーザレベルスレッドはカーネルレベルスレッドの集合上に
    多重化され，他のユーザレベルスレッドはカーネルレベルスレッドを専有
    する．カーネルレベルスレッドに恒久的に対応付けられたユーザレベル
    スレッドを\term[ばうんどすれっど]{バウンドスレッド}[bound thread]と
    呼ぶ．バウンドスレッドには，例えば高い優先度やよりよい応答性を
    与えられる．プロセスは複数のカーネルレベルスレッドを持つので，I/O で
    ブロックしたユーザレベルスレッドは自分のカーネルレベルスレッドだけを
    ブロックし，他のユーザレベルスレッドは残りのカーネルレベルスレッド上
    で実行を続ける．このモデルは他の2つの利点を兼ね備えるが，その代償と
    してユーザレベルとカーネルレベルのスレッド管理機構の間の協調が必要に
    なる．
\end{description}


\subsection{マルチスレッドのスコープ}

スレッド管理は2つのレベルで行われ，そのレベルによって資源を共有する単位
が決まる．
\term[ぷろせすすこーぷ]{プロセススコープ}[process scope]では，ユーザ
レベルのスレッドライブラリが1つのプロセス内のスレッドを管理し，スレッド
はそのプロセスの他のスレッドとだけ資源を奪い合う．カーネルはスレッドを
個別には見ない．\term[しすてむすこーぷ]{システムスコープ}[system scope]%
では，スレッドはオペレーティングシステムのスレッド管理機構，特に CPU
スケジューラによってシステム全体で管理され，システム内のすべての
スレッドと資源を奪い合う．

\begin{example}
  プロセス P は6つのユーザレベルスレッドを，プロセス Q は2つを実行する．
  プロセススコープではカーネルには2つのプロセスしか見えず，それぞれに
  CPU を均等に割り当てると，P の各スレッドは CPU の
  $\frac{1}{2}\cdot\frac{1}{6} = \frac{1}{12}$ を，Q の各スレッドは
  $\frac{1}{4}$ を受け取る．システムスコープではカーネルに8つの
  スレッドすべてが見え，CPU をそれらの間で均等に分けると各スレッドは
  $\frac{1}{8}$ を受け取り，P 全体では $\frac{3}{4}$ を受け取る．
\end{example}

\section{マルチスレッドのパターン}
\label{sec:l03-patterns}

マルチスレッドアプリケーションは，一般にボス・ワーカー，パイプライン，
階層化の3つのパターンのいずれかに従って構成される．\source{P2L2，
マルチスレッドのパターン；スレッドプールについては Silberschatz ら
ch.~4．}

\subsection{ボス・ワーカー}

\term[ぼすわーかーぱたーん]{ボス・ワーカーパターン}[boss-workers pattern]%
では，1つのボススレッドが要求を受け付けてワーカースレッドに割り当て，
各ワーカーはタスク全体を実行する．すべての要求がボスを通るので，システム
のスループットはボスによって制限される．ボスが要求ごとに時間
$t_{\mathrm{boss}}$ を費やすなら，ワーカーがいくつあってもシステムは単位
時間当たり高々 $1/t_{\mathrm{boss}}$ 個の要求しか受け付けられない．
したがって，ボスが要求ごとに行う仕事はできるだけ少なくしなければならない．

\begin{figure}[htbp]
  \centering
  % Boss-workers pattern with a shared work queue.
% Usage: % Boss-workers pattern with a shared work queue.
% Usage: % Boss-workers pattern with a shared work queue.
% Usage: \input{figures/l03-boss-workers}   (width ~106 mm)
\begin{tikzpicture}[x=1mm, y=1mm, font=\footnotesize\sffamily,
  box/.style={draw, rounded corners=2pt, fill=ln-box, align=center, inner sep=1pt},
  cell/.style={draw, fill=white, minimum width=3mm, minimum height=5mm, inner sep=0pt},
  >={Stealth[length=1.6mm]},
  lab/.style={font=\scriptsize\sffamily, text=ln-muted}]
  \node[lab, anchor=south west] at (0,0.5) {\iflangja{要求}{requests}};
  \node[box, minimum width=14mm, minimum height=7mm] (boss) at (26,0) {\iflangja{ボス}{boss}};
  \draw[->] (0,0) -- (boss);
  \foreach \i in {0,...,3} { \node[cell] at (46+3*\i,0) {}; }
  \node[lab, anchor=north] at (50.5,-3) {\iflangja{キュー}{queue}};
  \draw[->] (boss) -- (44.5,0);
  \foreach \y/\n in {8/1,0/2,-8/3} {
    \node[box, minimum width=18mm, minimum height=6mm] (w\n) at (82,\y) {\iflangja{ワーカー}{worker} \n};
    \draw[->] (56.5,0) -- (w\n.west);
    \draw[->] (w\n.east) -- ++(12,0);
  }
  \node[lab, anchor=south west] at (93,8.5) {\iflangja{応答}{responses}};
\end{tikzpicture}
   (width ~106 mm)
\begin{tikzpicture}[x=1mm, y=1mm, font=\footnotesize\sffamily,
  box/.style={draw, rounded corners=2pt, fill=ln-box, align=center, inner sep=1pt},
  cell/.style={draw, fill=white, minimum width=3mm, minimum height=5mm, inner sep=0pt},
  >={Stealth[length=1.6mm]},
  lab/.style={font=\scriptsize\sffamily, text=ln-muted}]
  \node[lab, anchor=south west] at (0,0.5) {\iflangja{要求}{requests}};
  \node[box, minimum width=14mm, minimum height=7mm] (boss) at (26,0) {\iflangja{ボス}{boss}};
  \draw[->] (0,0) -- (boss);
  \foreach \i in {0,...,3} { \node[cell] at (46+3*\i,0) {}; }
  \node[lab, anchor=north] at (50.5,-3) {\iflangja{キュー}{queue}};
  \draw[->] (boss) -- (44.5,0);
  \foreach \y/\n in {8/1,0/2,-8/3} {
    \node[box, minimum width=18mm, minimum height=6mm] (w\n) at (82,\y) {\iflangja{ワーカー}{worker} \n};
    \draw[->] (56.5,0) -- (w\n.west);
    \draw[->] (w\n.east) -- ++(12,0);
  }
  \node[lab, anchor=south west] at (93,8.5) {\iflangja{応答}{responses}};
\end{tikzpicture}
   (width ~106 mm)
\begin{tikzpicture}[x=1mm, y=1mm, font=\footnotesize\sffamily,
  box/.style={draw, rounded corners=2pt, fill=ln-box, align=center, inner sep=1pt},
  cell/.style={draw, fill=white, minimum width=3mm, minimum height=5mm, inner sep=0pt},
  >={Stealth[length=1.6mm]},
  lab/.style={font=\scriptsize\sffamily, text=ln-muted}]
  \node[lab, anchor=south west] at (0,0.5) {\iflangja{要求}{requests}};
  \node[box, minimum width=14mm, minimum height=7mm] (boss) at (26,0) {\iflangja{ボス}{boss}};
  \draw[->] (0,0) -- (boss);
  \foreach \i in {0,...,3} { \node[cell] at (46+3*\i,0) {}; }
  \node[lab, anchor=north] at (50.5,-3) {\iflangja{キュー}{queue}};
  \draw[->] (boss) -- (44.5,0);
  \foreach \y/\n in {8/1,0/2,-8/3} {
    \node[box, minimum width=18mm, minimum height=6mm] (w\n) at (82,\y) {\iflangja{ワーカー}{worker} \n};
    \draw[->] (56.5,0) -- (w\n.west);
    \draw[->] (w\n.east) -- ++(12,0);
  }
  \node[lab, anchor=south west] at (93,8.5) {\iflangja{応答}{responses}};
\end{tikzpicture}

  \caption{作業キューを用いるボス・ワーカー．ボスは要求をキューに入れる
    だけであり，各ワーカーはキューから要求を取り出してタスク全体を実行
    する．}
  \label{fig:l03-boss-workers}
\end{figure}

ボスがワーカーに仕事を渡す方法は2つある（\cref{fig:l03-boss-workers} は
2番目を示す）．
\begin{itemize}
  \item \textbf{特定のワーカーに直接通知する．}ワーカー間の同期は不要
    だが，ボスは各ワーカーが何をしているか，どれが空いているかを把握
    しなければならない．これは $t_{\mathrm{boss}}$ を増やし，スループット
    を下げる．
  \item \textbf{仕事をキューに置く．}ワーカーはそこから仕事を取り出す．
    これは生産者・消費者の構成である．ボスはワーカーについて何も知る必要
    がないが，ワーカー（とボス）は共有キューへのアクセスを同期しなければ
    ならない．ボスの仕事を少なく保てるので，通常はこちらの方がスループット
    がよい．
\end{itemize}

ワーカーは要求が到着したときに必要に応じて生成することもできるが，要求
ごとにスレッドを生成するとレイテンシが増える．より一般的には，
アプリケーションは事前に生成したワーカーの集合である
\term[すれっどぷーる]{スレッドプール}[thread pool]を用いる．プールは静的
でも動的でもよい．動的なプールは，キューが長くなると（典型的には1つずつ
ではなく一度に複数のスレッドずつ）拡大し，ワーカーがアイドルになると
縮小する．

このパターンは単純だが，スレッドプールの管理が必要であり，局所性を無視
する．ボスは，どのワーカーが似たタスクを直前に実行し，次のタスクを
ホットキャッシュで実行できるかを知らない．そこで変種として，すべての
ワーカーを同等にする代わりに，特定の種類の要求にワーカーを特化させるもの
がある．特化したワーカーは局所性がよく，アプリケーションはより重要な
種類の要求により多くのスレッドを割り当てることで異なるサービス品質を
提供できる．代償は負荷分散であり，各種類にいくつのスレッドを割り当てる
べきかの判断が難しくなる．

\subsection{パイプライン}

\term[ぱいぷらいんぱたーん]{パイプラインパターン}[pipeline pattern]では，
タスクを一連のサブタスクに分割し，各サブタスクをそれ専用のスレッド（または
スレッドのグループ）が実行する．ステージは共有バッファを通じて通信するので，
複数の要求がそれぞれ異なるステージで並行に処理される
（\cref{fig:l03-pipeline}）．スループットは最も弱い環，すなわち最も遅い
ステージで決まる．対策は，各ステージをスレッドプールにし，遅いステージ
により多くのスレッドを与えることである．ステージ $i$ が要求当たり時間
$t_i$ を要し，$p_i$ 個のスレッドを持つとき，スループットは
\begin{equation}
  X \;=\; \min_i \frac{p_i}{t_i}
  \label{eq:l03-pipeline-throughput}
\end{equation}
である．

\begin{figure}[htbp]
  \centering
  % Pipeline pattern: stages connected by buffers; stage 2 is a thread pool.
% Usage: % Pipeline pattern: stages connected by buffers; stage 2 is a thread pool.
% Usage: % Pipeline pattern: stages connected by buffers; stage 2 is a thread pool.
% Usage: \input{figures/l03-pipeline}   (width ~116 mm)
\begin{tikzpicture}[x=1mm, y=1mm, font=\footnotesize\sffamily,
  box/.style={draw, rounded corners=2pt, fill=ln-box, align=center, inner sep=1pt},
  th/.style={draw, circle, fill=white, minimum size=4mm, inner sep=0pt},
  cell/.style={draw, fill=white, minimum width=3mm, minimum height=5mm, inner sep=0pt},
  >={Stealth[length=1.6mm]},
  lab/.style={font=\scriptsize\sffamily, text=ln-muted}]
  \node[box, minimum width=14mm, minimum height=17mm] (s1) at (20,0) {};
  \node[box, minimum width=14mm, minimum height=17mm] (s2) at (60,0) {};
  \node[box, minimum width=14mm, minimum height=17mm] (s3) at (100,0) {};
  \foreach \s/\n in {s1/1,s2/2,s3/3} {
    \node[font=\scriptsize\sffamily, anchor=south] at (\s.north) {\iflangja{ステージ}{stage} \n};
  }
  \node[th] at (20,0) {};
  \foreach \y in {5,0,-5} { \node[th] at (60,\y) {}; }
  \node[th] at (100,0) {};
  \draw[->] (0,0) -- (s1);
  \foreach \i in {0,...,3} { \node[cell] at (35.5+3*\i,0) {}; }
  \foreach \i in {0,...,3} { \node[cell] at (75.5+3*\i,0) {}; }
  \draw[->] (s1) -- (34,0);  \draw[->] (45,0) -- (s2);
  \draw[->] (s2) -- (74,0);  \draw[->] (85,0) -- (s3);
  \draw[->] (s3) -- (116,0);
  \node[lab, anchor=north] at (40,-3) {\iflangja{バッファ}{buffer}};
  \node[lab, anchor=north] at (80,-3) {\iflangja{バッファ}{buffer}};
  \node[lab, anchor=north] at (60,-9) {\iflangja{スレッドプール}{thread pool}};
\end{tikzpicture}
   (width ~116 mm)
\begin{tikzpicture}[x=1mm, y=1mm, font=\footnotesize\sffamily,
  box/.style={draw, rounded corners=2pt, fill=ln-box, align=center, inner sep=1pt},
  th/.style={draw, circle, fill=white, minimum size=4mm, inner sep=0pt},
  cell/.style={draw, fill=white, minimum width=3mm, minimum height=5mm, inner sep=0pt},
  >={Stealth[length=1.6mm]},
  lab/.style={font=\scriptsize\sffamily, text=ln-muted}]
  \node[box, minimum width=14mm, minimum height=17mm] (s1) at (20,0) {};
  \node[box, minimum width=14mm, minimum height=17mm] (s2) at (60,0) {};
  \node[box, minimum width=14mm, minimum height=17mm] (s3) at (100,0) {};
  \foreach \s/\n in {s1/1,s2/2,s3/3} {
    \node[font=\scriptsize\sffamily, anchor=south] at (\s.north) {\iflangja{ステージ}{stage} \n};
  }
  \node[th] at (20,0) {};
  \foreach \y in {5,0,-5} { \node[th] at (60,\y) {}; }
  \node[th] at (100,0) {};
  \draw[->] (0,0) -- (s1);
  \foreach \i in {0,...,3} { \node[cell] at (35.5+3*\i,0) {}; }
  \foreach \i in {0,...,3} { \node[cell] at (75.5+3*\i,0) {}; }
  \draw[->] (s1) -- (34,0);  \draw[->] (45,0) -- (s2);
  \draw[->] (s2) -- (74,0);  \draw[->] (85,0) -- (s3);
  \draw[->] (s3) -- (116,0);
  \node[lab, anchor=north] at (40,-3) {\iflangja{バッファ}{buffer}};
  \node[lab, anchor=north] at (80,-3) {\iflangja{バッファ}{buffer}};
  \node[lab, anchor=north] at (60,-9) {\iflangja{スレッドプール}{thread pool}};
\end{tikzpicture}
   (width ~116 mm)
\begin{tikzpicture}[x=1mm, y=1mm, font=\footnotesize\sffamily,
  box/.style={draw, rounded corners=2pt, fill=ln-box, align=center, inner sep=1pt},
  th/.style={draw, circle, fill=white, minimum size=4mm, inner sep=0pt},
  cell/.style={draw, fill=white, minimum width=3mm, minimum height=5mm, inner sep=0pt},
  >={Stealth[length=1.6mm]},
  lab/.style={font=\scriptsize\sffamily, text=ln-muted}]
  \node[box, minimum width=14mm, minimum height=17mm] (s1) at (20,0) {};
  \node[box, minimum width=14mm, minimum height=17mm] (s2) at (60,0) {};
  \node[box, minimum width=14mm, minimum height=17mm] (s3) at (100,0) {};
  \foreach \s/\n in {s1/1,s2/2,s3/3} {
    \node[font=\scriptsize\sffamily, anchor=south] at (\s.north) {\iflangja{ステージ}{stage} \n};
  }
  \node[th] at (20,0) {};
  \foreach \y in {5,0,-5} { \node[th] at (60,\y) {}; }
  \node[th] at (100,0) {};
  \draw[->] (0,0) -- (s1);
  \foreach \i in {0,...,3} { \node[cell] at (35.5+3*\i,0) {}; }
  \foreach \i in {0,...,3} { \node[cell] at (75.5+3*\i,0) {}; }
  \draw[->] (s1) -- (34,0);  \draw[->] (45,0) -- (s2);
  \draw[->] (s2) -- (74,0);  \draw[->] (85,0) -- (s3);
  \draw[->] (s3) -- (116,0);
  \node[lab, anchor=north] at (40,-3) {\iflangja{バッファ}{buffer}};
  \node[lab, anchor=north] at (80,-3) {\iflangja{バッファ}{buffer}};
  \node[lab, anchor=north] at (60,-9) {\iflangja{スレッドプール}{thread pool}};
\end{tikzpicture}

  \caption{3ステージのパイプライン．ステージはバッファを通じて要求を受け
    渡す．遅い第2ステージは3スレッドのプールである．}
  \label{fig:l03-pipeline}
\end{figure}
パイプラインでは各スレッドが同じサブタスクを繰り返し実行するので，特化と
局所性が得られる．コストは，ステージのバランスを保つ手間とバッファの同期
オーバーヘッドである．

\begin{example}
  3ステージのパイプラインのステージ時間が \SI{10}{\milli\second}，
  \SI{40}{\milli\second}，\SI{20}{\milli\second} であるとする．各ステージ
  1スレッドの場合，\cref{eq:l03-pipeline-throughput} から
  $\min(100, 25, 50) = 25$ \mbox{要求/秒}となり，ステージ~2 が制約となる．
  ステージ~2 を4スレッドにするとスループットは
  $\min(100, 100, 50) = 50$ \mbox{要求/秒}になり，今度はステージ~3 が
  ボトルネックとなる．
\end{example}

\subsection{階層化}

\term[かいそうかぱたーん]{階層化パターン}[layered pattern]では，関連する
サブタスクを層にまとめて各層を専用のスレッドが担当し，すべての要求が
すべての層を下って再び上っていく（\cref{fig:l03-layered}）．パイプライン
と同様に特化をもたらすが，粒度はより粗い．すべてのアプリケーションに適して
いるわけではなく，各層とその上下の層との間の同期が必要である．

\begin{figure}[htbp]
  \centering
  % Layered pattern: a request passes down through the layers and back up.
% Usage: % Layered pattern: a request passes down through the layers and back up.
% Usage: % Layered pattern: a request passes down through the layers and back up.
% Usage: \input{figures/l03-layered}   (width ~84 mm)
\begin{tikzpicture}[x=1mm, y=1mm, font=\footnotesize\sffamily,
  box/.style={draw, rounded corners=2pt, fill=ln-box, align=center, inner sep=1pt},
  th/.style={draw, circle, fill=white, minimum size=4mm, inner sep=0pt},
  >={Stealth[length=1.6mm]},
  lab/.style={font=\scriptsize\sffamily, text=ln-muted}]
  \foreach \y/\n in {0/1,-8/2,-16/3} {
    \node[box, minimum width=52mm, minimum height=6.5mm] at (42,\y) {};
    \node[font=\scriptsize\sffamily, anchor=west] at (18,\y) {\iflangja{層}{layer} \n};
    \foreach \x in {44,54,64} { \node[th] at (\x,\y) {}; }
  }
  \draw[->] (12,5) -- (12,-23);
  \draw[->] (12,-23) -- (72,-23) -- (72,5);
  \node[lab, anchor=east] at (11,-8) {\iflangja{要求}{request}};
  \node[lab, anchor=west] at (73,-8) {\iflangja{応答}{response}};
\end{tikzpicture}
   (width ~84 mm)
\begin{tikzpicture}[x=1mm, y=1mm, font=\footnotesize\sffamily,
  box/.style={draw, rounded corners=2pt, fill=ln-box, align=center, inner sep=1pt},
  th/.style={draw, circle, fill=white, minimum size=4mm, inner sep=0pt},
  >={Stealth[length=1.6mm]},
  lab/.style={font=\scriptsize\sffamily, text=ln-muted}]
  \foreach \y/\n in {0/1,-8/2,-16/3} {
    \node[box, minimum width=52mm, minimum height=6.5mm] at (42,\y) {};
    \node[font=\scriptsize\sffamily, anchor=west] at (18,\y) {\iflangja{層}{layer} \n};
    \foreach \x in {44,54,64} { \node[th] at (\x,\y) {}; }
  }
  \draw[->] (12,5) -- (12,-23);
  \draw[->] (12,-23) -- (72,-23) -- (72,5);
  \node[lab, anchor=east] at (11,-8) {\iflangja{要求}{request}};
  \node[lab, anchor=west] at (73,-8) {\iflangja{応答}{response}};
\end{tikzpicture}
   (width ~84 mm)
\begin{tikzpicture}[x=1mm, y=1mm, font=\footnotesize\sffamily,
  box/.style={draw, rounded corners=2pt, fill=ln-box, align=center, inner sep=1pt},
  th/.style={draw, circle, fill=white, minimum size=4mm, inner sep=0pt},
  >={Stealth[length=1.6mm]},
  lab/.style={font=\scriptsize\sffamily, text=ln-muted}]
  \foreach \y/\n in {0/1,-8/2,-16/3} {
    \node[box, minimum width=52mm, minimum height=6.5mm] at (42,\y) {};
    \node[font=\scriptsize\sffamily, anchor=west] at (18,\y) {\iflangja{層}{layer} \n};
    \foreach \x in {44,54,64} { \node[th] at (\x,\y) {}; }
  }
  \draw[->] (12,5) -- (12,-23);
  \draw[->] (12,-23) -- (72,-23) -- (72,5);
  \node[lab, anchor=east] at (11,-8) {\iflangja{要求}{request}};
  \node[lab, anchor=west] at (73,-8) {\iflangja{応答}{response}};
\end{tikzpicture}

  \caption{階層化パターン．各層は関連するサブタスクをまとめ，専用の
    スレッドを持つ．要求は層を下って再び上っていく．}
  \label{fig:l03-layered}
\end{figure}

\subsection{ボス・ワーカーとパイプラインの比較}

一群の要求をどちらのパターンが速く処理し終えるかは，要求の数に依存する．それ
ぞれが同じ総処理時間を要する $n$ 個の要求を考える．$w$ 個のワーカーを持ち
ボスの時間 $t_{\mathrm{boss}}$ が無視できるボス・ワーカー設計は，要求を
$w$ 個ずつ完了し，各グループにタスク全体の時間 $t$ がかかる．時間 $s$ の
等しい $k$ 個のステージからなるパイプラインは，$k$ ステージ時間の後に
最初の要求を完了し，その後は1ステージ時間ごとに1要求を完了する．したがって
\begin{equation}
  T_{\mathrm{bw}}(n) = \Bigl\lceil \frac{n}{w} \Bigr\rceil\, t ,
  \qquad
  T_{\mathrm{pl}}(n) = (k + n - 1)\, s
  \label{eq:l03-batch-time}
\end{equation}
である．

\begin{example}
  どちらの設計も4つのスレッドを使い，1要求に \SI{100}{\milli\second} の
  作業が必要であるとする．ボス・ワーカー設計は1つのボスと $w=3$ 個の
  ワーカーを持ち $t=\SI{100}{\milli\second}$，パイプラインは
  $s=\SI{25}{\milli\second}$ のステージを $k=4$ 個持つ．
  \SI{100}{\milli\second} 当たりに完了する要求はボス・ワーカーが3，定常
  状態のパイプラインが4なので，ボス・ワーカーが先行できるのは，パイプ
  ラインの立ち上がり時間 $k\,s$ がまだ効いている少数の要求においてだけで
  ある．そこでも結果は $n$ に依存する．\cref{eq:l03-batch-time} の
  $T_{\mathrm{bw}}$ は階段状に増えるからである．$n = 2$，3，6 では
  ボス・ワーカーが速く，$n = 4$，7，8 ではパイプラインが速く，
  $n = 1$，5，9 では両者が等しく，$n \geq 10$ では常にパイプラインが
  勝つ．
  \begin{center}
    \small
    \begin{tabular}{@{}rrr@{}}
      \toprule
      $n$ & ボス・ワーカー（\si{\milli\second}） & パイプライン（\si{\milli\second}） \\
      \midrule
       3 & 100 & 150 \\
       4 & 200 & 175 \\
       6 & 200 & 225 \\
      12 & 400 & 375 \\
      \bottomrule
    \end{tabular}
  \end{center}
\end{example}

\begin{keypoint}[章のまとめ]
  スレッドはプロセスのアドレス空間を共有する実行コンテキストであり，
  並列性，特化，アイドル時間の利用（\cref{eq:l03-idle}）をプロセスより
  安価に提供する．共有データはミューテックスで保護しなければならず，
  スレッドは条件変数で条件を待ち，その条件は常に \texttt{while} ループで
  再検査する．代理変数を用いるとリーダ・ライタのようなより複雑な方針を
  構築でき，
  ミューテックスの外で signal するとスプリアスウェイクアップを避けられ，
  一貫したロック順序はデッドロックを防ぐ．ユーザレベルスレッドは一対一，
  多対一，多対多のいずれかのモデルに従ってカーネルレベルスレッド上で実行
  され，アプリケーションはボス・ワーカー，パイプライン，階層として構成
  される．
\end{keypoint}

\end{document}
